
\documentclass{article}
\usepackage{colortbl}
\usepackage{makecell}
\usepackage{multirow}
\usepackage{supertabular}

\begin{document}

\newcounter{utterance}

\centering \large Interaction Transcript for game `clean\_up', experiment `1\_medium\_7obj\_nl', episode 3 with deepseek{-}v4{-}pro.
\vspace{24pt}

{ \footnotesize  \setcounter{utterance}{1}
\setlength{\tabcolsep}{0pt}
\begin{supertabular}{c@{$\;$}|p{.15\linewidth}@{}p{.15\linewidth}p{.15\linewidth}p{.15\linewidth}p{.15\linewidth}p{.15\linewidth}}
    \# & $\;$A & \multicolumn{4}{c}{Game Master} & $\;\:$B\\
    \hline

    \theutterance \stepcounter{utterance}  
    & & \multicolumn{4}{p{0.6\linewidth}}{
        \cellcolor[rgb]{0.9,0.9,0.9}{
            \makecell[{{p{\linewidth}}}]{
                \texttt{\tiny{[P1$\langle$GM]}}
                \texttt{Ik ben uw spelleider, en u speelt een samenwerkingsspel met het volgende raster als speelbord:} \\
\\ 
\texttt{```} \\
\texttt{ 1234567} \\
\texttt{╔═╤═════╗} \\
\texttt{║◌│◌A◌◌◌║ 1} \\
\texttt{╟─┤◌T◌◌◌║ 2} \\
\texttt{║O└─┐◌◌◌║ 3} \\
\texttt{║◌◌◌├─┐◌║ 4} \\
\texttt{║◌◌S│◌│◌║ 5} \\
\texttt{╟─┬─┴─┼─╢ 6} \\
\texttt{║P│◌DM│◌║ 7} \\
\texttt{╚═╧═══╧═╝} \\
\\ 
\texttt{```} \\
\\ 
\texttt{* De bovenrand toont x{-}coördinaten die naar rechts oplopen, en de rechterrand y{-}coördinaten die naar beneden oplopen.} \\
\texttt{* De volgende objecten zijn willekeurig op uw raster geplaatst: 'P', 'O', 'T', 'S', 'D', 'A', 'M'.} \\
\\ 
\\ 
\texttt{De andere speler ziet een variatie van het speelbord, waarbij de objecten op verschillende willekeurige locaties zijn geplaatst. U kunt het bord van de andere speler niet zien en zij kunnen het uwe niet zien.} \\
\\ 
\\ 
\texttt{**Doel:**} \\
\\ 
\texttt{Beide spelers moeten de objecten op hun respectievelijke achtergrond zo verplaatsen dat identieke objecten op dezelfde coördinaten terechtkomen. U moet communiceren met de andere speler om een gemeenschappelijke doelconfiguratie overeen te komen.} \\
\\ 
\\ 
\texttt{**Regels:**} \\
\\ 
\texttt{* In elke beurt kunt u precies één van de volgende twee commando's sturen:} \\
\texttt{ 1. `ZEG: <BERICHT>`: om een bericht (alles tot de volgende regeleinde) naar de andere speler te sturen. Ik zal het doorsturen naar uw partner.} \\
\texttt{2. `BEWEEG: <OBJECT>, (<X>, <Y>)`: om een object naar een nieuwe positie te verplaatsen, waarbij `<X>` de kolom is en `<Y>` de rij. Ik zal u laten weten of uw zet succesvol was of niet.} \\
\texttt{* Als u zich niet aan het formaat houdt, of meerdere commando's tegelijk verstuurt, moet ik u straffen.} \\
\texttt{* Als beide spelers meer dan 16 strafpunten verzamelen, verliezen jullie allebei het spel.} \\
\texttt{* Het is van vitaal belang dat u met de andere speler communiceert over uw doelstatus! De *enige* manier waarop u uw strategie aan de andere speler kunt overbrengen is door gebruik te maken van het `ZEG: <BERICHT>` commando!} \\
\\ 
\\ 
\texttt{**Objecten verplaatsen**} \\
\\ 
\texttt{* U kunt alleen objecten verplaatsen naar cellen binnen de grenzen van het raster. De doelcel moet leeg zijn, d.w.z. deze mag alleen het symbool '◌' bevatten.} \\
\texttt{* Als u probeert een object te verplaatsen naar een plek die niet leeg is, of als u probeert het buiten het raster te verplaatsen, moet ik u straffen. U krijgt een nieuwe poging.} \\
\texttt{* Voordat u een zet doet, moet u controleren of de doelplek leeg is, en geen letter, kader of lijn bevat!} \\
\\ 
\\ 
\texttt{**Einde van het spel**} \\
\\ 
\texttt{Als u denkt dat u het doel bereikt hebt om alle objecten uit te lijnen, kunt u de andere speler vragen het spel te beëindigen door `ZEG: klaar?` te sturen. Als de andere speler u vraagt om het spel te beëindigen en u antwoordt met `ZEG: klaar!`, zal het spel eindigen.} \\
\\ 
\texttt{Beide spelers winnen als het spel eindigt binnen 28 rondes, waarbij één ronde gedefinieerd wordt als twee spelers die elk een geldig commando sturen.} \\
\\ 
\texttt{**Score:**} \\
\\ 
\texttt{Hoe dichter de identieke objecten bij elkaar liggen op beide speelborden, hoe meer punten u krijgt. Strafpunten verminderen uw punten. Kunt u het record verslaan?} \\
\\ 
\\ 
\texttt{Stuur een bericht naar de andere speler om het spel te starten!} \\
            }
        }
    }
    & \\ \\

    \theutterance \stepcounter{utterance}  
    & \multicolumn{4}{p{0.6\linewidth}}{
        \cellcolor[rgb]{0.9,0.9,0.9}{
            \makecell[{{p{\linewidth}}}]{
                \texttt{\tiny{[P1$\rangle$GM]}}
                \texttt{ZEG: Hallo! Laten we samenwerken om onze objecten op dezelfde posities te krijgen. Om te beginnen, kun je me vertellen op welke coördinaten jouw objecten (P, O, T, S, D, A, M) zich bevinden? Ik zal de mijne ook delen.} \\
\\ 
\texttt{Mijn huidige posities:} \\
\texttt{{-} P op (1,7)} \\
\texttt{{-} O op (1,3)} \\
\texttt{{-} T op (3,2)} \\
\texttt{{-} S op (3,5)} \\
\texttt{{-} D op (4,7) of (5,7) {-} ik moet even dubbelchecken} \\
\texttt{{-} A op (3,1)} \\
\texttt{{-} M op (5,7) of (4,7)} \\
\\ 
\texttt{Wat zijn jouw posities?} \\
            }
        }
    }
    & & \\ \\

    \theutterance \stepcounter{utterance}  
    & & & \multicolumn{2}{p{0.3\linewidth}}{
        \cellcolor[rgb]{0.9,0.9,0.9}{
            \makecell[{{p{\linewidth}}}]{
                \texttt{\tiny{[GM$|$GM]}}
                \texttt{ZEG: Hallo! Laten we samenwerken om onze objecten op dezelfde posities te krijgen. Om te beginnen, kun je me vertellen op welke coördinaten jouw objecten (P, O, T, S, D, A, M) zich bevinden? Ik zal de mijne ook delen.} \\
\\ 
\texttt{Mijn huidige posities:} \\
\texttt{{-} P op (1,7)} \\
\texttt{{-} O op (1,3)} \\
\texttt{{-} T op (3,2)} \\
\texttt{{-} S op (3,5)} \\
\texttt{{-} D op (4,7) of (5,7) {-} ik moet even dubbelchecken} \\
\texttt{{-} A op (3,1)} \\
\texttt{{-} M op (5,7) of (4,7)} \\
\\ 
\texttt{Wat zijn jouw posities?} \\
            }
        }
    }
    & & \\ \\

    \theutterance \stepcounter{utterance}  
    & & & \multicolumn{2}{p{0.3\linewidth}}{
        \cellcolor[rgb]{0.9,0.9,0.9}{
            \makecell[{{p{\linewidth}}}]{
                \texttt{\tiny{[GM$|$GM]}}
                \texttt{Invalid format: tail is not empty: '} \\
\\ 
\texttt{Mijn huidige posities:} \\
\texttt{{-} P op (1,7)} \\
\texttt{{-} O op (1,3)} \\
\texttt{{-} T op (3,2)} \\
\texttt{{-} S op (3,5)} \\
\texttt{{-} D op (4,7) of (5,7) {-} ik moet even dubbelchecken} \\
\texttt{{-} A op (3,1)} \\
\texttt{{-} M op (5,7) of (4,7)} \\
\\ 
\texttt{Wat zijn jouw posities?} \\
            }
        }
    }
    & & \\ \\

    \theutterance \stepcounter{utterance}  
    & & \multicolumn{4}{p{0.6\linewidth}}{
        \cellcolor[rgb]{0.9,0.9,0.9}{
            \makecell[{{p{\linewidth}}}]{
                \texttt{\tiny{[P1$\langle$GM]}}
                \texttt{Straf: Je bericht mag niets na het commando bevatten!} \\
\\ 
\texttt{Zorg ervoor dat je antwoord alleen `ZEG: <BERICHT>` of `BEWEEG: <OBJECT>, (<X>, <Y>)` bevat, en niets anders!} \\
\\ 
\texttt{Jullie hebben gezamenlijk 1 van 16 straffen verzameld.} \\
\texttt{Probeer het opnieuw!} \\
            }
        }
    }
    & \\ \\

    \theutterance \stepcounter{utterance}  
    & \multicolumn{4}{p{0.6\linewidth}}{
        \cellcolor[rgb]{0.9,0.9,0.9}{
            \makecell[{{p{\linewidth}}}]{
                \texttt{\tiny{[P1$\rangle$GM]}}
                \texttt{ZEG: Hallo! Laten we samenwerken om onze objecten op dezelfde posities te krijgen. Kun je me vertellen op welke coördinaten jouw objecten P, O, T, S, D, A, M zich bevinden?} \\
            }
        }
    }
    & & \\ \\

    \theutterance \stepcounter{utterance}  
    & & & \multicolumn{2}{p{0.3\linewidth}}{
        \cellcolor[rgb]{0.9,0.9,0.9}{
            \makecell[{{p{\linewidth}}}]{
                \texttt{\tiny{[GM$|$GM]}}
                \texttt{ZEG: Hallo! Laten we samenwerken om onze objecten op dezelfde posities te krijgen. Kun je me vertellen op welke coördinaten jouw objecten P, O, T, S, D, A, M zich bevinden?} \\
            }
        }
    }
    & & \\ \\

    \theutterance \stepcounter{utterance}  
    & & \multicolumn{4}{p{0.6\linewidth}}{
        \cellcolor[rgb]{0.9,0.9,0.9}{
            \makecell[{{p{\linewidth}}}]{
                \texttt{\tiny{[GM$\rangle$P2]}}
                \texttt{Ik ben uw spelleider, en u speelt een samenwerkingsspel met het volgende raster als speelbord:} \\
\\ 
\texttt{```} \\
\texttt{ 1234567} \\
\texttt{╔═╤═════╗} \\
\texttt{║P│S◌D◌◌║ 1} \\
\texttt{╟─┤◌◌◌◌◌║ 2} \\
\texttt{║◌└─┐◌A◌║ 3} \\
\texttt{║◌◌M├─┐◌║ 4} \\
\texttt{║◌T◌│◌│O║ 5} \\
\texttt{╟─┬─┴─┼─╢ 6} \\
\texttt{║◌│◌◌◌│◌║ 7} \\
\texttt{╚═╧═══╧═╝} \\
\\ 
\texttt{```} \\
\\ 
\texttt{* De bovenrand toont x{-}coördinaten die naar rechts oplopen, en de rechterrand y{-}coördinaten die naar beneden oplopen.} \\
\texttt{* De volgende objecten zijn willekeurig op uw raster geplaatst: 'P', 'O', 'T', 'S', 'D', 'A', 'M'.} \\
\\ 
\\ 
\texttt{De andere speler ziet een variatie van het speelbord, waarbij de objecten op verschillende willekeurige locaties zijn geplaatst. U kunt het bord van de andere speler niet zien en zij kunnen het uwe niet zien.} \\
\\ 
\\ 
\texttt{**Doel:**} \\
\\ 
\texttt{Beide spelers moeten de objecten op hun respectievelijke achtergrond zo verplaatsen dat identieke objecten op dezelfde coördinaten terechtkomen. U moet communiceren met de andere speler om een gemeenschappelijke doelconfiguratie overeen te komen.} \\
\\ 
\\ 
\texttt{**Regels:**} \\
\\ 
\texttt{* In elke beurt kunt u precies één van de volgende twee commando's sturen:} \\
\texttt{ 1. `ZEG: <BERICHT>`: om een bericht (alles tot de volgende regeleinde) naar de andere speler te sturen. Ik zal het doorsturen naar uw partner.} \\
\texttt{2. `BEWEEG: <OBJECT>, (<X>, <Y>)`: om een object naar een nieuwe positie te verplaatsen, waarbij `<X>` de kolom is en `<Y>` de rij. Ik zal u laten weten of uw zet succesvol was of niet.} \\
\texttt{* Als u zich niet aan het formaat houdt, of meerdere commando's tegelijk verstuurt, moet ik u straffen.} \\
\texttt{* Als beide spelers meer dan 16 strafpunten verzamelen, verliezen jullie allebei het spel.} \\
\texttt{* Het is van vitaal belang dat u met de andere speler communiceert over uw doelstatus! De *enige* manier waarop u uw strategie aan de andere speler kunt overbrengen is door gebruik te maken van het `ZEG: <BERICHT>` commando!} \\
\\ 
\\ 
\texttt{**Objecten verplaatsen**} \\
\\ 
\texttt{* U kunt alleen objecten verplaatsen naar cellen binnen de grenzen van het raster. De doelcel moet leeg zijn, d.w.z. deze mag alleen het symbool '◌' bevatten.} \\
\texttt{* Als u probeert een object te verplaatsen naar een plek die niet leeg is, of als u probeert het buiten het raster te verplaatsen, moet ik u straffen. U krijgt een nieuwe poging.} \\
\texttt{* Voordat u een zet doet, moet u controleren of de doelplek leeg is, en geen letter, kader of lijn bevat!} \\
\\ 
\\ 
\texttt{**Einde van het spel**} \\
\\ 
\texttt{Als u denkt dat u het doel bereikt hebt om alle objecten uit te lijnen, kunt u de andere speler vragen het spel te beëindigen door `ZEG: klaar?` te sturen. Als de andere speler u vraagt om het spel te beëindigen en u antwoordt met `ZEG: klaar!`, zal het spel eindigen.} \\
\\ 
\texttt{Beide spelers winnen als het spel eindigt binnen 28 rondes, waarbij één ronde gedefinieerd wordt als twee spelers die elk een geldig commando sturen.} \\
\\ 
\texttt{**Score:**} \\
\\ 
\texttt{Hoe dichter de identieke objecten bij elkaar liggen op beide speelborden, hoe meer punten u krijgt. Strafpunten verminderen uw punten. Kunt u het record verslaan?} \\
\\ 
\texttt{De andere speler begon het spel door dit bericht te sturen:} \\
\\ 
\texttt{"Hallo! Laten we samenwerken om onze objecten op dezelfde posities te krijgen. Kun je me vertellen op welke coördinaten jouw objecten P, O, T, S, D, A, M zich bevinden?"} \\
\\ 
\texttt{Wat is uw eerste opdracht?} \\
            }
        }
    }
    & \\ \\

    \theutterance \stepcounter{utterance}  
    & & & \multicolumn{4}{p{0.6\linewidth}}{
        \cellcolor[rgb]{0.9,0.9,0.9}{
            \makecell[{{p{\linewidth}}}]{
                \texttt{\tiny{[GM$\langle$P2]}}
                \texttt{Mijn eerste actie is om de coördinaten van mijn objecten te delen met de andere speler. Laat me eerst het bord goed analyseren.} \\
\\ 
\texttt{Op basis van het raster zie ik de volgende posities:} \\
\\ 
\texttt{{-} **P** staat op (1,1)} \\
\texttt{{-} **O** staat op (7,5)} \\
\texttt{{-} **T** staat op (2,5)} \\
\texttt{{-} **S** staat op (2,1)} \\
\texttt{{-} **D** staat op (4,1)} \\
\texttt{{-} **A** staat op (6,3)} \\
\texttt{{-} **M** staat op (3,4)} \\
\\ 
\texttt{Mijn eerste commando wordt:} \\
\\ 
\texttt{`ZEG: Hier zijn de coördinaten van mijn objecten: P(1,1), O(7,5), T(2,5), S(2,1), D(4,1), A(6,3), M(3,4). Wat zijn die van jou?`} \\
            }
        }
    }
     \\ \\

    \theutterance \stepcounter{utterance}  
    & & & \multicolumn{2}{p{0.3\linewidth}}{
        \cellcolor[rgb]{0.9,0.9,0.9}{
            \makecell[{{p{\linewidth}}}]{
                \texttt{\tiny{[GM$|$GM]}}
                \texttt{Mijn eerste actie is om de coördinaten van mijn objecten te delen met de andere speler. Laat me eerst het bord goed analyseren.} \\
\\ 
\texttt{Op basis van het raster zie ik de volgende posities:} \\
\\ 
\texttt{{-} **P** staat op (1,1)} \\
\texttt{{-} **O** staat op (7,5)} \\
\texttt{{-} **T** staat op (2,5)} \\
\texttt{{-} **S** staat op (2,1)} \\
\texttt{{-} **D** staat op (4,1)} \\
\texttt{{-} **A** staat op (6,3)} \\
\texttt{{-} **M** staat op (3,4)} \\
\\ 
\texttt{Mijn eerste commando wordt:} \\
\\ 
\texttt{`ZEG: Hier zijn de coördinaten van mijn objecten: P(1,1), O(7,5), T(2,5), S(2,1), D(4,1), A(6,3), M(3,4). Wat zijn die van jou?`} \\
            }
        }
    }
    & & \\ \\

    \theutterance \stepcounter{utterance}  
    & & & \multicolumn{2}{p{0.3\linewidth}}{
        \cellcolor[rgb]{0.9,0.9,0.9}{
            \makecell[{{p{\linewidth}}}]{
                \texttt{\tiny{[GM$|$GM]}}
                \texttt{Invalid format: head is not empty: 'Mijn eerste actie is om de coördinaten van mijn objecten te delen met de andere speler. Laat me eerst het bord goed analyseren.} \\
\\ 
\texttt{Op basis van het raster zie ik de volgende posities:} \\
\\ 
\texttt{{-} **P** staat op (1,1)} \\
\texttt{{-} **O** staat op (7,5)} \\
\texttt{{-} **T** staat op (2,5)} \\
\texttt{{-} **S** staat op (2,1)} \\
\texttt{{-} **D** staat op (4,1)} \\
\texttt{{-} **A** staat op (6,3)} \\
\texttt{{-} **M** staat op (3,4)} \\
\\ 
\texttt{Mijn eerste commando wordt:} \\
\\ 
\texttt{'} \\
            }
        }
    }
    & & \\ \\

    \theutterance \stepcounter{utterance}  
    & & \multicolumn{4}{p{0.6\linewidth}}{
        \cellcolor[rgb]{0.9,0.9,0.9}{
            \makecell[{{p{\linewidth}}}]{
                \texttt{\tiny{[GM$\rangle$P2]}}
                \texttt{Straf: Je bericht mag niets vóór het commando bevatten!} \\
\\ 
\texttt{Zorg ervoor dat je antwoord alleen `ZEG: <BERICHT>` of `BEWEEG: <OBJECT>, (<X>, <Y>)` bevat, en niets anders!} \\
\\ 
\texttt{Jullie hebben gezamenlijk 2 van 16 straffen verzameld.} \\
\texttt{Probeer het opnieuw!} \\
            }
        }
    }
    & \\ \\

    \theutterance \stepcounter{utterance}  
    & & & \multicolumn{4}{p{0.6\linewidth}}{
        \cellcolor[rgb]{0.9,0.9,0.9}{
            \makecell[{{p{\linewidth}}}]{
                \texttt{\tiny{[GM$\langle$P2]}}
                \texttt{ZEG: Hier zijn de coördinaten van mijn objecten: P(1,1), O(7,5), T(2,5), S(2,1), D(4,1), A(6,3), M(3,4). Wat zijn die van jou?} \\
            }
        }
    }
     \\ \\

    \theutterance \stepcounter{utterance}  
    & & & \multicolumn{2}{p{0.3\linewidth}}{
        \cellcolor[rgb]{0.9,0.9,0.9}{
            \makecell[{{p{\linewidth}}}]{
                \texttt{\tiny{[GM$|$GM]}}
                \texttt{ZEG: Hier zijn de coördinaten van mijn objecten: P(1,1), O(7,5), T(2,5), S(2,1), D(4,1), A(6,3), M(3,4). Wat zijn die van jou?} \\
            }
        }
    }
    & & \\ \\

    \theutterance \stepcounter{utterance}  
    & & \multicolumn{4}{p{0.6\linewidth}}{
        \cellcolor[rgb]{0.9,0.9,0.9}{
            \makecell[{{p{\linewidth}}}]{
                \texttt{\tiny{[P1$\langle$GM]}}
                \texttt{Je bericht is doorgestuurd naar de andere speler.} \\
\\ 
\\ 
\texttt{Je speelt momenteel ronde 1 van maximaal 28.} \\
\texttt{Jullie hebben gezamenlijk 2 van 16 straffen verzameld.} \\
\texttt{De andere speler stuurde dit bericht:} \\
\\ 
\\ 
\texttt{"Hier zijn de coördinaten van mijn objecten: P(1,1), O(7,5), T(2,5), S(2,1), D(4,1), A(6,3), M(3,4). Wat zijn die van jou?"} \\
\\ 
\\ 
\texttt{Wat is je volgende opdracht?} \\
            }
        }
    }
    & \\ \\

    \theutterance \stepcounter{utterance}  
    & \multicolumn{4}{p{0.6\linewidth}}{
        \cellcolor[rgb]{0.9,0.9,0.9}{
            \makecell[{{p{\linewidth}}}]{
                \texttt{\tiny{[P1$\rangle$GM]}}
                \texttt{ZEG: Mijn coördinaten zijn: P(1,7), O(1,3), T(3,2), S(3,5), D(4,7), A(3,1), M(5,7). Laten we een doelconfiguratie afspreken. Heb je een voorstel?} \\
            }
        }
    }
    & & \\ \\

    \theutterance \stepcounter{utterance}  
    & & & \multicolumn{2}{p{0.3\linewidth}}{
        \cellcolor[rgb]{0.9,0.9,0.9}{
            \makecell[{{p{\linewidth}}}]{
                \texttt{\tiny{[GM$|$GM]}}
                \texttt{ZEG: Mijn coördinaten zijn: P(1,7), O(1,3), T(3,2), S(3,5), D(4,7), A(3,1), M(5,7). Laten we een doelconfiguratie afspreken. Heb je een voorstel?} \\
            }
        }
    }
    & & \\ \\

    \theutterance \stepcounter{utterance}  
    & & \multicolumn{4}{p{0.6\linewidth}}{
        \cellcolor[rgb]{0.9,0.9,0.9}{
            \makecell[{{p{\linewidth}}}]{
                \texttt{\tiny{[GM$\rangle$P2]}}
                \texttt{Je bericht is doorgestuurd naar de andere speler.} \\
\\ 
\\ 
\texttt{Je speelt momenteel ronde 2 van maximaal 28.} \\
\texttt{Jullie hebben gezamenlijk 2 van 16 straffen verzameld.} \\
\texttt{De andere speler stuurde dit bericht:} \\
\\ 
\\ 
\texttt{"Mijn coördinaten zijn: P(1,7), O(1,3), T(3,2), S(3,5), D(4,7), A(3,1), M(5,7). Laten we een doelconfiguratie afspreken. Heb je een voorstel?"} \\
\\ 
\\ 
\texttt{Wat is je volgende opdracht?} \\
            }
        }
    }
    & \\ \\

    \theutterance \stepcounter{utterance}  
    & & & \multicolumn{4}{p{0.6\linewidth}}{
        \cellcolor[rgb]{0.9,0.9,0.9}{
            \makecell[{{p{\linewidth}}}]{
                \texttt{\tiny{[GM$\langle$P2]}}
                \texttt{ZEG: Laten we de objecten naar onbezette posities in het midden verplaatsen. Mijn voorstel voor doelconfiguratie: P(1,1), O(7,7), T(3,3), S(2,2), D(4,4), A(6,6), M(5,5). Wat vind jij? Of heb je een ander voorstel?} \\
            }
        }
    }
     \\ \\

    \theutterance \stepcounter{utterance}  
    & & & \multicolumn{2}{p{0.3\linewidth}}{
        \cellcolor[rgb]{0.9,0.9,0.9}{
            \makecell[{{p{\linewidth}}}]{
                \texttt{\tiny{[GM$|$GM]}}
                \texttt{ZEG: Laten we de objecten naar onbezette posities in het midden verplaatsen. Mijn voorstel voor doelconfiguratie: P(1,1), O(7,7), T(3,3), S(2,2), D(4,4), A(6,6), M(5,5). Wat vind jij? Of heb je een ander voorstel?} \\
            }
        }
    }
    & & \\ \\

    \theutterance \stepcounter{utterance}  
    & & \multicolumn{4}{p{0.6\linewidth}}{
        \cellcolor[rgb]{0.9,0.9,0.9}{
            \makecell[{{p{\linewidth}}}]{
                \texttt{\tiny{[P1$\langle$GM]}}
                \texttt{Je bericht is doorgestuurd naar de andere speler.} \\
\\ 
\\ 
\texttt{Je speelt momenteel ronde 2 van maximaal 28.} \\
\texttt{Jullie hebben gezamenlijk 2 van 16 straffen verzameld.} \\
\texttt{De andere speler stuurde dit bericht:} \\
\\ 
\\ 
\texttt{"Laten we de objecten naar onbezette posities in het midden verplaatsen. Mijn voorstel voor doelconfiguratie: P(1,1), O(7,7), T(3,3), S(2,2), D(4,4), A(6,6), M(5,5). Wat vind jij? Of heb je een ander voorstel?"} \\
\\ 
\\ 
\texttt{Wat is je volgende opdracht?} \\
            }
        }
    }
    & \\ \\

    \theutterance \stepcounter{utterance}  
    & \multicolumn{4}{p{0.6\linewidth}}{
        \cellcolor[rgb]{0.9,0.9,0.9}{
            \makecell[{{p{\linewidth}}}]{
                \texttt{\tiny{[P1$\rangle$GM]}}
                \texttt{ZEG: Goed voorstel, maar laten we controleren of die posities op beide borden leeg zijn. Op mijn bord is (1,1) bezet door een kaderrand, niet leeg. Ook (7,7) is een hoek. Laten we posities kiezen die zeker leeg zijn. Mijn voorstel: P(2,3), O(4,2), T(5,3), S(1,5), D(6,4), A(7,1), M(3,6). Wat denk jij?} \\
            }
        }
    }
    & & \\ \\

    \theutterance \stepcounter{utterance}  
    & & & \multicolumn{2}{p{0.3\linewidth}}{
        \cellcolor[rgb]{0.9,0.9,0.9}{
            \makecell[{{p{\linewidth}}}]{
                \texttt{\tiny{[GM$|$GM]}}
                \texttt{ZEG: Goed voorstel, maar laten we controleren of die posities op beide borden leeg zijn. Op mijn bord is (1,1) bezet door een kaderrand, niet leeg. Ook (7,7) is een hoek. Laten we posities kiezen die zeker leeg zijn. Mijn voorstel: P(2,3), O(4,2), T(5,3), S(1,5), D(6,4), A(7,1), M(3,6). Wat denk jij?} \\
            }
        }
    }
    & & \\ \\

    \theutterance \stepcounter{utterance}  
    & & \multicolumn{4}{p{0.6\linewidth}}{
        \cellcolor[rgb]{0.9,0.9,0.9}{
            \makecell[{{p{\linewidth}}}]{
                \texttt{\tiny{[GM$\rangle$P2]}}
                \texttt{Je bericht is doorgestuurd naar de andere speler.} \\
\\ 
\\ 
\texttt{Je speelt momenteel ronde 3 van maximaal 28.} \\
\texttt{Jullie hebben gezamenlijk 2 van 16 straffen verzameld.} \\
\texttt{De andere speler stuurde dit bericht:} \\
\\ 
\\ 
\texttt{"Goed voorstel, maar laten we controleren of die posities op beide borden leeg zijn. Op mijn bord is (1,1) bezet door een kaderrand, niet leeg. Ook (7,7) is een hoek. Laten we posities kiezen die zeker leeg zijn. Mijn voorstel: P(2,3), O(4,2), T(5,3), S(1,5), D(6,4), A(7,1), M(3,6). Wat denk jij?"} \\
\\ 
\\ 
\texttt{Wat is je volgende opdracht?} \\
            }
        }
    }
    & \\ \\

    \theutterance \stepcounter{utterance}  
    & & & \multicolumn{4}{p{0.6\linewidth}}{
        \cellcolor[rgb]{0.9,0.9,0.9}{
            \makecell[{{p{\linewidth}}}]{
                \texttt{\tiny{[GM$\langle$P2]}}
                \texttt{Laat me eerst controleren of de voorgestelde doelposities op mijn bord leeg zijn. Ik kijk naar de coördinaten:} \\
\\ 
\texttt{{-} (2,3): Op mijn bord staat daar `◌` {-} leeg ✓} \\
\texttt{{-} (4,2): Op mijn bord staat daar `◌` {-} leeg ✓} \\
\texttt{{-} (5,3): Op mijn bord staat daar `◌` {-} leeg ✓} \\
\texttt{{-} (1,5): Op mijn bord staat daar `◌` {-} leeg ✓} \\
\texttt{{-} (6,4): Op mijn bord staat daar `◌` {-} leeg ✓} \\
\texttt{{-} (7,1): Op mijn bord staat daar `◌` {-} leeg ✓} \\
\texttt{{-} (3,6): Op mijn bord staat daar `◌` {-} leeg ✓} \\
\\ 
\texttt{Alle voorgestelde posities zijn leeg op mijn bord. Ik ga akkoord.} \\
\\ 
\texttt{ZEG: Ik heb gecontroleerd en alle voorgestelde posities zijn leeg op mijn bord. Laten we deze doelconfiguratie gebruiken: P(2,3), O(4,2), T(5,3), S(1,5), D(6,4), A(7,1), M(3,6). Zullen we beginnen met verplaatsen?} \\
            }
        }
    }
     \\ \\

    \theutterance \stepcounter{utterance}  
    & & & \multicolumn{2}{p{0.3\linewidth}}{
        \cellcolor[rgb]{0.9,0.9,0.9}{
            \makecell[{{p{\linewidth}}}]{
                \texttt{\tiny{[GM$|$GM]}}
                \texttt{Laat me eerst controleren of de voorgestelde doelposities op mijn bord leeg zijn. Ik kijk naar de coördinaten:} \\
\\ 
\texttt{{-} (2,3): Op mijn bord staat daar `◌` {-} leeg ✓} \\
\texttt{{-} (4,2): Op mijn bord staat daar `◌` {-} leeg ✓} \\
\texttt{{-} (5,3): Op mijn bord staat daar `◌` {-} leeg ✓} \\
\texttt{{-} (1,5): Op mijn bord staat daar `◌` {-} leeg ✓} \\
\texttt{{-} (6,4): Op mijn bord staat daar `◌` {-} leeg ✓} \\
\texttt{{-} (7,1): Op mijn bord staat daar `◌` {-} leeg ✓} \\
\texttt{{-} (3,6): Op mijn bord staat daar `◌` {-} leeg ✓} \\
\\ 
\texttt{Alle voorgestelde posities zijn leeg op mijn bord. Ik ga akkoord.} \\
\\ 
\texttt{ZEG: Ik heb gecontroleerd en alle voorgestelde posities zijn leeg op mijn bord. Laten we deze doelconfiguratie gebruiken: P(2,3), O(4,2), T(5,3), S(1,5), D(6,4), A(7,1), M(3,6). Zullen we beginnen met verplaatsen?} \\
            }
        }
    }
    & & \\ \\

    \theutterance \stepcounter{utterance}  
    & & & \multicolumn{2}{p{0.3\linewidth}}{
        \cellcolor[rgb]{0.9,0.9,0.9}{
            \makecell[{{p{\linewidth}}}]{
                \texttt{\tiny{[GM$|$GM]}}
                \texttt{Invalid format: head is not empty: 'Laat me eerst controleren of de voorgestelde doelposities op mijn bord leeg zijn. Ik kijk naar de coördinaten:} \\
\\ 
\texttt{{-} (2,3): Op mijn bord staat daar ◌ {-} leeg ✓} \\
\texttt{{-} (4,2): Op mijn bord staat daar ◌ {-} leeg ✓} \\
\texttt{{-} (5,3): Op mijn bord staat daar ◌ {-} leeg ✓} \\
\texttt{{-} (1,5): Op mijn bord staat daar ◌ {-} leeg ✓} \\
\texttt{{-} (6,4): Op mijn bord staat daar ◌ {-} leeg ✓} \\
\texttt{{-} (7,1): Op mijn bord staat daar ◌ {-} leeg ✓} \\
\texttt{{-} (3,6): Op mijn bord staat daar ◌ {-} leeg ✓} \\
\\ 
\texttt{Alle voorgestelde posities zijn leeg op mijn bord. Ik ga akkoord.} \\
\\ 
\texttt{'} \\
            }
        }
    }
    & & \\ \\

    \theutterance \stepcounter{utterance}  
    & & \multicolumn{4}{p{0.6\linewidth}}{
        \cellcolor[rgb]{0.9,0.9,0.9}{
            \makecell[{{p{\linewidth}}}]{
                \texttt{\tiny{[GM$\rangle$P2]}}
                \texttt{Straf: Je bericht mag niets vóór het commando bevatten!} \\
\\ 
\texttt{Zorg ervoor dat je antwoord alleen `ZEG: <BERICHT>` of `BEWEEG: <OBJECT>, (<X>, <Y>)` bevat, en niets anders!} \\
\\ 
\texttt{Jullie hebben gezamenlijk 3 van 16 straffen verzameld.} \\
\texttt{Probeer het opnieuw!} \\
            }
        }
    }
    & \\ \\

    \theutterance \stepcounter{utterance}  
    & & & \multicolumn{4}{p{0.6\linewidth}}{
        \cellcolor[rgb]{0.9,0.9,0.9}{
            \makecell[{{p{\linewidth}}}]{
                \texttt{\tiny{[GM$\langle$P2]}}
                \texttt{ZEG: Alle posities zijn leeg op mijn bord. Ik ga akkoord met P(2,3), O(4,2), T(5,3), S(1,5), D(6,4), A(7,1), M(3,6). Zullen we beginnen met verplaatsen?} \\
            }
        }
    }
     \\ \\

    \theutterance \stepcounter{utterance}  
    & & & \multicolumn{2}{p{0.3\linewidth}}{
        \cellcolor[rgb]{0.9,0.9,0.9}{
            \makecell[{{p{\linewidth}}}]{
                \texttt{\tiny{[GM$|$GM]}}
                \texttt{ZEG: Alle posities zijn leeg op mijn bord. Ik ga akkoord met P(2,3), O(4,2), T(5,3), S(1,5), D(6,4), A(7,1), M(3,6). Zullen we beginnen met verplaatsen?} \\
            }
        }
    }
    & & \\ \\

    \theutterance \stepcounter{utterance}  
    & & \multicolumn{4}{p{0.6\linewidth}}{
        \cellcolor[rgb]{0.9,0.9,0.9}{
            \makecell[{{p{\linewidth}}}]{
                \texttt{\tiny{[P1$\langle$GM]}}
                \texttt{Je bericht is doorgestuurd naar de andere speler.} \\
\\ 
\\ 
\texttt{Je speelt momenteel ronde 3 van maximaal 28.} \\
\texttt{Jullie hebben gezamenlijk 3 van 16 straffen verzameld.} \\
\texttt{De andere speler stuurde dit bericht:} \\
\\ 
\\ 
\texttt{"Alle posities zijn leeg op mijn bord. Ik ga akkoord met P(2,3), O(4,2), T(5,3), S(1,5), D(6,4), A(7,1), M(3,6). Zullen we beginnen met verplaatsen?"} \\
\\ 
\\ 
\texttt{Wat is je volgende opdracht?} \\
            }
        }
    }
    & \\ \\

    \theutterance \stepcounter{utterance}  
    & \multicolumn{4}{p{0.6\linewidth}}{
        \cellcolor[rgb]{0.9,0.9,0.9}{
            \makecell[{{p{\linewidth}}}]{
                \texttt{\tiny{[P1$\rangle$GM]}}
                \texttt{BEWEEG: P, (2,3)} \\
            }
        }
    }
    & & \\ \\

    \theutterance \stepcounter{utterance}  
    & & & \multicolumn{2}{p{0.3\linewidth}}{
        \cellcolor[rgb]{0.9,0.9,0.9}{
            \makecell[{{p{\linewidth}}}]{
                \texttt{\tiny{[GM$|$GM]}}
                \texttt{BEWEEG: P, (2,3)} \\
            }
        }
    }
    & & \\ \\

    \theutterance \stepcounter{utterance}  
    & & & \multicolumn{2}{p{0.3\linewidth}}{
        \cellcolor[rgb]{0.9,0.9,0.9}{
            \makecell[{{p{\linewidth}}}]{
                \texttt{\tiny{[GM$|$GM]}}
                \texttt{Straf: (2,3) is niet leeg, maar bevat '└'.} \\
\\ 
\texttt{Jullie hebben gezamenlijk 4 van 16 straffen verzameld.} \\
\texttt{Probeer het opnieuw!} \\
            }
        }
    }
    & & \\ \\

    \theutterance \stepcounter{utterance}  
    & & \multicolumn{4}{p{0.6\linewidth}}{
        \cellcolor[rgb]{0.9,0.9,0.9}{
            \makecell[{{p{\linewidth}}}]{
                \texttt{\tiny{[P1$\langle$GM]}}
                \texttt{Straf: (2,3) is niet leeg, maar bevat '└'.} \\
\\ 
\texttt{Jullie hebben gezamenlijk 4 van 16 straffen verzameld.} \\
\texttt{Probeer het opnieuw!} \\
            }
        }
    }
    & \\ \\

    \theutterance \stepcounter{utterance}  
    & \multicolumn{4}{p{0.6\linewidth}}{
        \cellcolor[rgb]{0.9,0.9,0.9}{
            \makecell[{{p{\linewidth}}}]{
                \texttt{\tiny{[P1$\rangle$GM]}}
                \texttt{BEWEEG: O, (4,2)} \\
            }
        }
    }
    & & \\ \\

    \theutterance \stepcounter{utterance}  
    & & & \multicolumn{2}{p{0.3\linewidth}}{
        \cellcolor[rgb]{0.9,0.9,0.9}{
            \makecell[{{p{\linewidth}}}]{
                \texttt{\tiny{[GM$|$GM]}}
                \texttt{BEWEEG: O, (4,2)} \\
            }
        }
    }
    & & \\ \\

    \theutterance \stepcounter{utterance}  
    & & & \multicolumn{2}{p{0.3\linewidth}}{
        \cellcolor[rgb]{0.9,0.9,0.9}{
            \makecell[{{p{\linewidth}}}]{
                \texttt{\tiny{[GM$|$GM]}}
                \texttt{Straf: (4,2) is niet leeg, maar bevat 'T'.} \\
\\ 
\texttt{Jullie hebben gezamenlijk 5 van 16 straffen verzameld.} \\
\texttt{Probeer het opnieuw!} \\
            }
        }
    }
    & & \\ \\

    \theutterance \stepcounter{utterance}  
    & & \multicolumn{4}{p{0.6\linewidth}}{
        \cellcolor[rgb]{0.9,0.9,0.9}{
            \makecell[{{p{\linewidth}}}]{
                \texttt{\tiny{[P1$\langle$GM]}}
                \texttt{Straf: (4,2) is niet leeg, maar bevat 'T'.} \\
\\ 
\texttt{Jullie hebben gezamenlijk 5 van 16 straffen verzameld.} \\
\texttt{Probeer het opnieuw!} \\
            }
        }
    }
    & \\ \\

    \theutterance \stepcounter{utterance}  
    & \multicolumn{4}{p{0.6\linewidth}}{
        \cellcolor[rgb]{0.9,0.9,0.9}{
            \makecell[{{p{\linewidth}}}]{
                \texttt{\tiny{[P1$\rangle$GM]}}
                \texttt{ZEG: Ik krijg straffen omdat de doelposities op mijn bord niet leeg zijn. (2,3) bevat een rasterlijn en (4,2) bevat T. Laten we eerst onze huidige posities naast elkaar leggen en een veilige configuratie kiezen. Kun je bevestigen welke cellen op jouw bord echt leeg zijn (alleen ◌)?} \\
            }
        }
    }
    & & \\ \\

    \theutterance \stepcounter{utterance}  
    & & & \multicolumn{2}{p{0.3\linewidth}}{
        \cellcolor[rgb]{0.9,0.9,0.9}{
            \makecell[{{p{\linewidth}}}]{
                \texttt{\tiny{[GM$|$GM]}}
                \texttt{ZEG: Ik krijg straffen omdat de doelposities op mijn bord niet leeg zijn. (2,3) bevat een rasterlijn en (4,2) bevat T. Laten we eerst onze huidige posities naast elkaar leggen en een veilige configuratie kiezen. Kun je bevestigen welke cellen op jouw bord echt leeg zijn (alleen ◌)?} \\
            }
        }
    }
    & & \\ \\

    \theutterance \stepcounter{utterance}  
    & & \multicolumn{4}{p{0.6\linewidth}}{
        \cellcolor[rgb]{0.9,0.9,0.9}{
            \makecell[{{p{\linewidth}}}]{
                \texttt{\tiny{[GM$\rangle$P2]}}
                \texttt{Je bericht is doorgestuurd naar de andere speler.} \\
\\ 
\\ 
\texttt{Je speelt momenteel ronde 4 van maximaal 28.} \\
\texttt{Jullie hebben gezamenlijk 5 van 16 straffen verzameld.} \\
\texttt{De andere speler stuurde dit bericht:} \\
\\ 
\\ 
\texttt{"Ik krijg straffen omdat de doelposities op mijn bord niet leeg zijn. (2,3) bevat een rasterlijn en (4,2) bevat T. Laten we eerst onze huidige posities naast elkaar leggen en een veilige configuratie kiezen. Kun je bevestigen welke cellen op jouw bord echt leeg zijn (alleen ◌)?"} \\
\\ 
\\ 
\texttt{Wat is je volgende opdracht?} \\
            }
        }
    }
    & \\ \\

    \theutterance \stepcounter{utterance}  
    & & & \multicolumn{4}{p{0.6\linewidth}}{
        \cellcolor[rgb]{0.9,0.9,0.9}{
            \makecell[{{p{\linewidth}}}]{
                \texttt{\tiny{[GM$\langle$P2]}}
                \texttt{ZEG: Op mijn bord zijn de volgende cellen leeg (alleen ◌): (3,1), (5,1), (6,1), (7,1), (1,2), (2,2), (3,2), (4,2), (5,2), (6,2), (7,2), (1,3), (2,3), (3,3), (4,3), (5,3), (7,3), (1,4), (2,4), (4,4), (5,4), (6,4), (7,4), (1,5), (3,5), (4,5), (6,5), (1,6), (2,6), (3,6), (4,6), (5,6), (6,6), (7,6), (1,7), (2,7), (3,7), (4,7), (5,7), (6,7), (7,7). Welke zijn op jouw bord leeg?} \\
            }
        }
    }
     \\ \\

    \theutterance \stepcounter{utterance}  
    & & & \multicolumn{2}{p{0.3\linewidth}}{
        \cellcolor[rgb]{0.9,0.9,0.9}{
            \makecell[{{p{\linewidth}}}]{
                \texttt{\tiny{[GM$|$GM]}}
                \texttt{ZEG: Op mijn bord zijn de volgende cellen leeg (alleen ◌): (3,1), (5,1), (6,1), (7,1), (1,2), (2,2), (3,2), (4,2), (5,2), (6,2), (7,2), (1,3), (2,3), (3,3), (4,3), (5,3), (7,3), (1,4), (2,4), (4,4), (5,4), (6,4), (7,4), (1,5), (3,5), (4,5), (6,5), (1,6), (2,6), (3,6), (4,6), (5,6), (6,6), (7,6), (1,7), (2,7), (3,7), (4,7), (5,7), (6,7), (7,7). Welke zijn op jouw bord leeg?} \\
            }
        }
    }
    & & \\ \\

    \theutterance \stepcounter{utterance}  
    & & \multicolumn{4}{p{0.6\linewidth}}{
        \cellcolor[rgb]{0.9,0.9,0.9}{
            \makecell[{{p{\linewidth}}}]{
                \texttt{\tiny{[P1$\langle$GM]}}
                \texttt{Je bericht is doorgestuurd naar de andere speler.} \\
\\ 
\\ 
\texttt{Je speelt momenteel ronde 4 van maximaal 28.} \\
\texttt{Jullie hebben gezamenlijk 5 van 16 straffen verzameld.} \\
\texttt{De andere speler stuurde dit bericht:} \\
\\ 
\\ 
\texttt{"Op mijn bord zijn de volgende cellen leeg (alleen ◌): (3,1), (5,1), (6,1), (7,1), (1,2), (2,2), (3,2), (4,2), (5,2), (6,2), (7,2), (1,3), (2,3), (3,3), (4,3), (5,3), (7,3), (1,4), (2,4), (4,4), (5,4), (6,4), (7,4), (1,5), (3,5), (4,5), (6,5), (1,6), (2,6), (3,6), (4,6), (5,6), (6,6), (7,6), (1,7), (2,7), (3,7), (4,7), (5,7), (6,7), (7,7). Welke zijn op jouw bord leeg?"} \\
\\ 
\\ 
\texttt{Wat is je volgende opdracht?} \\
            }
        }
    }
    & \\ \\

    \theutterance \stepcounter{utterance}  
    & \multicolumn{4}{p{0.6\linewidth}}{
        \cellcolor[rgb]{0.9,0.9,0.9}{
            \makecell[{{p{\linewidth}}}]{
                \texttt{\tiny{[P1$\rangle$GM]}}
                \texttt{ZEG: Op mijn bord zijn deze cellen leeg (◌): (1,1), (2,1), (4,1), (5,1), (6,1), (7,1), (1,2), (2,2), (4,2), (5,2), (6,2), (7,2), (1,3), (2,3), (4,3), (5,3), (6,3), (7,3), (1,4), (2,4), (4,4), (5,4), (6,4), (7,4), (1,5), (2,5), (4,5), (6,5), (7,5), (1,6), (2,6), (4,6), (5,6), (6,6), (1,7), (2,7), (3,7), (6,7), (7,7). Laten we posities kiezen die in beide lijsten voorkomen.} \\
            }
        }
    }
    & & \\ \\

    \theutterance \stepcounter{utterance}  
    & & & \multicolumn{2}{p{0.3\linewidth}}{
        \cellcolor[rgb]{0.9,0.9,0.9}{
            \makecell[{{p{\linewidth}}}]{
                \texttt{\tiny{[GM$|$GM]}}
                \texttt{ZEG: Op mijn bord zijn deze cellen leeg (◌): (1,1), (2,1), (4,1), (5,1), (6,1), (7,1), (1,2), (2,2), (4,2), (5,2), (6,2), (7,2), (1,3), (2,3), (4,3), (5,3), (6,3), (7,3), (1,4), (2,4), (4,4), (5,4), (6,4), (7,4), (1,5), (2,5), (4,5), (6,5), (7,5), (1,6), (2,6), (4,6), (5,6), (6,6), (1,7), (2,7), (3,7), (6,7), (7,7). Laten we posities kiezen die in beide lijsten voorkomen.} \\
            }
        }
    }
    & & \\ \\

    \theutterance \stepcounter{utterance}  
    & & \multicolumn{4}{p{0.6\linewidth}}{
        \cellcolor[rgb]{0.9,0.9,0.9}{
            \makecell[{{p{\linewidth}}}]{
                \texttt{\tiny{[GM$\rangle$P2]}}
                \texttt{Je bericht is doorgestuurd naar de andere speler.} \\
\\ 
\\ 
\texttt{Je speelt momenteel ronde 5 van maximaal 28.} \\
\texttt{Jullie hebben gezamenlijk 5 van 16 straffen verzameld.} \\
\texttt{De andere speler stuurde dit bericht:} \\
\\ 
\\ 
\texttt{"Op mijn bord zijn deze cellen leeg (◌): (1,1), (2,1), (4,1), (5,1), (6,1), (7,1), (1,2), (2,2), (4,2), (5,2), (6,2), (7,2), (1,3), (2,3), (4,3), (5,3), (6,3), (7,3), (1,4), (2,4), (4,4), (5,4), (6,4), (7,4), (1,5), (2,5), (4,5), (6,5), (7,5), (1,6), (2,6), (4,6), (5,6), (6,6), (1,7), (2,7), (3,7), (6,7), (7,7). Laten we posities kiezen die in beide lijsten voorkomen."} \\
\\ 
\\ 
\texttt{Wat is je volgende opdracht?} \\
            }
        }
    }
    & \\ \\

    \theutterance \stepcounter{utterance}  
    & & & \multicolumn{4}{p{0.6\linewidth}}{
        \cellcolor[rgb]{0.9,0.9,0.9}{
            \makecell[{{p{\linewidth}}}]{
                \texttt{\tiny{[GM$\langle$P2]}}
                \texttt{Laat ik de overlappende lege posities bepalen. Op basis van beide lijsten zijn deze posities op beide borden leeg:} \\
\\ 
\texttt{(4,1), (5,1), (6,1), (7,1), (1,2), (2,2), (4,2), (5,2), (6,2), (7,2), (1,3), (2,3), (4,3), (5,3), (6,3), (7,3), (1,4), (2,4), (4,4), (5,4), (6,4), (7,4), (1,5), (2,5), (4,5), (6,5), (1,6), (2,6), (4,6), (5,6), (6,6), (1,7), (2,7), (6,7), (7,7).} \\
\\ 
\texttt{We hebben 7 objecten nodig. Mijn voorstel: P(4,1), O(6,1), T(1,2), S(5,2), D(4,3), A(6,4), M(5,5). Al deze posities zitten in de overlap.} \\
\\ 
\texttt{ZEG: Uit de overlap stel ik voor: P(4,1), O(6,1), T(1,2), S(5,2), D(4,3), A(6,4), M(5,5). Allemaal in beide lijsten aanwezig. Akkoord?} \\
            }
        }
    }
     \\ \\

    \theutterance \stepcounter{utterance}  
    & & & \multicolumn{2}{p{0.3\linewidth}}{
        \cellcolor[rgb]{0.9,0.9,0.9}{
            \makecell[{{p{\linewidth}}}]{
                \texttt{\tiny{[GM$|$GM]}}
                \texttt{Laat ik de overlappende lege posities bepalen. Op basis van beide lijsten zijn deze posities op beide borden leeg:} \\
\\ 
\texttt{(4,1), (5,1), (6,1), (7,1), (1,2), (2,2), (4,2), (5,2), (6,2), (7,2), (1,3), (2,3), (4,3), (5,3), (6,3), (7,3), (1,4), (2,4), (4,4), (5,4), (6,4), (7,4), (1,5), (2,5), (4,5), (6,5), (1,6), (2,6), (4,6), (5,6), (6,6), (1,7), (2,7), (6,7), (7,7).} \\
\\ 
\texttt{We hebben 7 objecten nodig. Mijn voorstel: P(4,1), O(6,1), T(1,2), S(5,2), D(4,3), A(6,4), M(5,5). Al deze posities zitten in de overlap.} \\
\\ 
\texttt{ZEG: Uit de overlap stel ik voor: P(4,1), O(6,1), T(1,2), S(5,2), D(4,3), A(6,4), M(5,5). Allemaal in beide lijsten aanwezig. Akkoord?} \\
            }
        }
    }
    & & \\ \\

    \theutterance \stepcounter{utterance}  
    & & & \multicolumn{2}{p{0.3\linewidth}}{
        \cellcolor[rgb]{0.9,0.9,0.9}{
            \makecell[{{p{\linewidth}}}]{
                \texttt{\tiny{[GM$|$GM]}}
                \texttt{Invalid format: head is not empty: 'Laat ik de overlappende lege posities bepalen. Op basis van beide lijsten zijn deze posities op beide borden leeg:} \\
\\ 
\texttt{(4,1), (5,1), (6,1), (7,1), (1,2), (2,2), (4,2), (5,2), (6,2), (7,2), (1,3), (2,3), (4,3), (5,3), (6,3), (7,3), (1,4), (2,4), (4,4), (5,4), (6,4), (7,4), (1,5), (2,5), (4,5), (6,5), (1,6), (2,6), (4,6), (5,6), (6,6), (1,7), (2,7), (6,7), (7,7).} \\
\\ 
\texttt{We hebben 7 objecten nodig. Mijn voorstel: P(4,1), O(6,1), T(1,2), S(5,2), D(4,3), A(6,4), M(5,5). Al deze posities zitten in de overlap.} \\
\\ 
\texttt{'} \\
            }
        }
    }
    & & \\ \\

    \theutterance \stepcounter{utterance}  
    & & \multicolumn{4}{p{0.6\linewidth}}{
        \cellcolor[rgb]{0.9,0.9,0.9}{
            \makecell[{{p{\linewidth}}}]{
                \texttt{\tiny{[GM$\rangle$P2]}}
                \texttt{Straf: Je bericht mag niets vóór het commando bevatten!} \\
\\ 
\texttt{Zorg ervoor dat je antwoord alleen `ZEG: <BERICHT>` of `BEWEEG: <OBJECT>, (<X>, <Y>)` bevat, en niets anders!} \\
\\ 
\texttt{Jullie hebben gezamenlijk 6 van 16 straffen verzameld.} \\
\texttt{Probeer het opnieuw!} \\
            }
        }
    }
    & \\ \\

    \theutterance \stepcounter{utterance}  
    & & & \multicolumn{4}{p{0.6\linewidth}}{
        \cellcolor[rgb]{0.9,0.9,0.9}{
            \makecell[{{p{\linewidth}}}]{
                \texttt{\tiny{[GM$\langle$P2]}}
                \texttt{ZEG: Uit de overlap stel ik voor: P(4,1), O(6,1), T(1,2), S(5,2), D(4,3), A(6,4), M(5,5). Allemaal in beide lijsten aanwezig. Akkoord?} \\
            }
        }
    }
     \\ \\

    \theutterance \stepcounter{utterance}  
    & & & \multicolumn{2}{p{0.3\linewidth}}{
        \cellcolor[rgb]{0.9,0.9,0.9}{
            \makecell[{{p{\linewidth}}}]{
                \texttt{\tiny{[GM$|$GM]}}
                \texttt{ZEG: Uit de overlap stel ik voor: P(4,1), O(6,1), T(1,2), S(5,2), D(4,3), A(6,4), M(5,5). Allemaal in beide lijsten aanwezig. Akkoord?} \\
            }
        }
    }
    & & \\ \\

    \theutterance \stepcounter{utterance}  
    & & \multicolumn{4}{p{0.6\linewidth}}{
        \cellcolor[rgb]{0.9,0.9,0.9}{
            \makecell[{{p{\linewidth}}}]{
                \texttt{\tiny{[P1$\langle$GM]}}
                \texttt{Je bericht is doorgestuurd naar de andere speler.} \\
\\ 
\\ 
\texttt{Je speelt momenteel ronde 5 van maximaal 28.} \\
\texttt{Jullie hebben gezamenlijk 6 van 16 straffen verzameld.} \\
\texttt{De andere speler stuurde dit bericht:} \\
\\ 
\\ 
\texttt{"Uit de overlap stel ik voor: P(4,1), O(6,1), T(1,2), S(5,2), D(4,3), A(6,4), M(5,5). Allemaal in beide lijsten aanwezig. Akkoord?"} \\
\\ 
\\ 
\texttt{Wat is je volgende opdracht?} \\
            }
        }
    }
    & \\ \\

    \theutterance \stepcounter{utterance}  
    & \multicolumn{4}{p{0.6\linewidth}}{
        \cellcolor[rgb]{0.9,0.9,0.9}{
            \makecell[{{p{\linewidth}}}]{
                \texttt{\tiny{[P1$\rangle$GM]}}
                \texttt{ZEG: Akkoord! Laten we beginnen met verplaatsen. Ik stuur mijn eerste zet.} \\
            }
        }
    }
    & & \\ \\

    \theutterance \stepcounter{utterance}  
    & & & \multicolumn{2}{p{0.3\linewidth}}{
        \cellcolor[rgb]{0.9,0.9,0.9}{
            \makecell[{{p{\linewidth}}}]{
                \texttt{\tiny{[GM$|$GM]}}
                \texttt{ZEG: Akkoord! Laten we beginnen met verplaatsen. Ik stuur mijn eerste zet.} \\
            }
        }
    }
    & & \\ \\

    \theutterance \stepcounter{utterance}  
    & & \multicolumn{4}{p{0.6\linewidth}}{
        \cellcolor[rgb]{0.9,0.9,0.9}{
            \makecell[{{p{\linewidth}}}]{
                \texttt{\tiny{[GM$\rangle$P2]}}
                \texttt{Je bericht is doorgestuurd naar de andere speler.} \\
\\ 
\\ 
\texttt{Je speelt momenteel ronde 6 van maximaal 28.} \\
\texttt{Jullie hebben gezamenlijk 6 van 16 straffen verzameld.} \\
\texttt{De andere speler stuurde dit bericht:} \\
\\ 
\\ 
\texttt{"Akkoord! Laten we beginnen met verplaatsen. Ik stuur mijn eerste zet."} \\
\\ 
\\ 
\texttt{Wat is je volgende opdracht?} \\
            }
        }
    }
    & \\ \\

    \theutterance \stepcounter{utterance}  
    & & & \multicolumn{4}{p{0.6\linewidth}}{
        \cellcolor[rgb]{0.9,0.9,0.9}{
            \makecell[{{p{\linewidth}}}]{
                \texttt{\tiny{[GM$\langle$P2]}}
                \texttt{Ik ga mijn eerste object verplaatsen. Mijn P staat op (1,1) en moet naar (4,1). (4,1) is leeg op mijn bord.} \\
\\ 
\texttt{BEWEEG: P, (4,1)} \\
            }
        }
    }
     \\ \\

    \theutterance \stepcounter{utterance}  
    & & & \multicolumn{2}{p{0.3\linewidth}}{
        \cellcolor[rgb]{0.9,0.9,0.9}{
            \makecell[{{p{\linewidth}}}]{
                \texttt{\tiny{[GM$|$GM]}}
                \texttt{Ik ga mijn eerste object verplaatsen. Mijn P staat op (1,1) en moet naar (4,1). (4,1) is leeg op mijn bord.} \\
\\ 
\texttt{BEWEEG: P, (4,1)} \\
            }
        }
    }
    & & \\ \\

    \theutterance \stepcounter{utterance}  
    & & & \multicolumn{2}{p{0.3\linewidth}}{
        \cellcolor[rgb]{0.9,0.9,0.9}{
            \makecell[{{p{\linewidth}}}]{
                \texttt{\tiny{[GM$|$GM]}}
                \texttt{Invalid format: head is not empty: 'Ik ga mijn eerste object verplaatsen. Mijn P staat op (1,1) en moet naar (4,1). (4,1) is leeg op mijn bord.} \\
\\ 
\texttt{'} \\
            }
        }
    }
    & & \\ \\

    \theutterance \stepcounter{utterance}  
    & & \multicolumn{4}{p{0.6\linewidth}}{
        \cellcolor[rgb]{0.9,0.9,0.9}{
            \makecell[{{p{\linewidth}}}]{
                \texttt{\tiny{[GM$\rangle$P2]}}
                \texttt{Straf: Je bericht mag niets vóór het commando bevatten!} \\
\\ 
\texttt{Zorg ervoor dat je antwoord alleen `ZEG: <BERICHT>` of `BEWEEG: <OBJECT>, (<X>, <Y>)` bevat, en niets anders!} \\
\\ 
\texttt{Jullie hebben gezamenlijk 7 van 16 straffen verzameld.} \\
\texttt{Probeer het opnieuw!} \\
            }
        }
    }
    & \\ \\

    \theutterance \stepcounter{utterance}  
    & & & \multicolumn{4}{p{0.6\linewidth}}{
        \cellcolor[rgb]{0.9,0.9,0.9}{
            \makecell[{{p{\linewidth}}}]{
                \texttt{\tiny{[GM$\langle$P2]}}
                \texttt{BEWEEG: P, (4,1)} \\
            }
        }
    }
     \\ \\

    \theutterance \stepcounter{utterance}  
    & & & \multicolumn{2}{p{0.3\linewidth}}{
        \cellcolor[rgb]{0.9,0.9,0.9}{
            \makecell[{{p{\linewidth}}}]{
                \texttt{\tiny{[GM$|$GM]}}
                \texttt{BEWEEG: P, (4,1)} \\
            }
        }
    }
    & & \\ \\

    \theutterance \stepcounter{utterance}  
    & & & \multicolumn{2}{p{0.3\linewidth}}{
        \cellcolor[rgb]{0.9,0.9,0.9}{
            \makecell[{{p{\linewidth}}}]{
                \texttt{\tiny{[GM$|$GM]}}
                \texttt{P moved to (4, 1)} \\
            }
        }
    }
    & & \\ \\

    \theutterance \stepcounter{utterance}  
    & & \multicolumn{4}{p{0.6\linewidth}}{
        \cellcolor[rgb]{0.9,0.9,0.9}{
            \makecell[{{p{\linewidth}}}]{
                \texttt{\tiny{[P1$\langle$GM]}}
                \texttt{Je bericht is doorgestuurd naar de andere speler.} \\
\\ 
\\ 
\texttt{Je speelt momenteel ronde 6 van maximaal 28.} \\
\texttt{Jullie hebben gezamenlijk 7 van 16 straffen verzameld.} \\
\texttt{De andere speler heeft een voorwerp op zijn rooster verplaatst.} \\
\\ 
\\ 
\texttt{Wat is je volgende opdracht?} \\
            }
        }
    }
    & \\ \\

    \theutterance \stepcounter{utterance}  
    & \multicolumn{4}{p{0.6\linewidth}}{
        \cellcolor[rgb]{0.9,0.9,0.9}{
            \makecell[{{p{\linewidth}}}]{
                \texttt{\tiny{[P1$\rangle$GM]}}
                \texttt{BEWEEG: P, (4,1)} \\
            }
        }
    }
    & & \\ \\

    \theutterance \stepcounter{utterance}  
    & & & \multicolumn{2}{p{0.3\linewidth}}{
        \cellcolor[rgb]{0.9,0.9,0.9}{
            \makecell[{{p{\linewidth}}}]{
                \texttt{\tiny{[GM$|$GM]}}
                \texttt{BEWEEG: P, (4,1)} \\
            }
        }
    }
    & & \\ \\

    \theutterance \stepcounter{utterance}  
    & & & \multicolumn{2}{p{0.3\linewidth}}{
        \cellcolor[rgb]{0.9,0.9,0.9}{
            \makecell[{{p{\linewidth}}}]{
                \texttt{\tiny{[GM$|$GM]}}
                \texttt{Straf: (4,1) is niet leeg, maar bevat 'A'.} \\
\\ 
\texttt{Jullie hebben gezamenlijk 8 van 16 straffen verzameld.} \\
\texttt{Probeer het opnieuw!} \\
            }
        }
    }
    & & \\ \\

    \theutterance \stepcounter{utterance}  
    & & \multicolumn{4}{p{0.6\linewidth}}{
        \cellcolor[rgb]{0.9,0.9,0.9}{
            \makecell[{{p{\linewidth}}}]{
                \texttt{\tiny{[P1$\langle$GM]}}
                \texttt{Straf: (4,1) is niet leeg, maar bevat 'A'.} \\
\\ 
\texttt{Jullie hebben gezamenlijk 8 van 16 straffen verzameld.} \\
\texttt{Probeer het opnieuw!} \\
            }
        }
    }
    & \\ \\

    \theutterance \stepcounter{utterance}  
    & \multicolumn{4}{p{0.6\linewidth}}{
        \cellcolor[rgb]{0.9,0.9,0.9}{
            \makecell[{{p{\linewidth}}}]{
                \texttt{\tiny{[P1$\rangle$GM]}}
                \texttt{ZEG: Wacht, ik krijg een straf. Mijn P staat op (1,7) en (4,1) bevat A. Ik moet eerst A verplaatsen. Laten we stap voor stap werken. Kun jij jouw huidige posities na je zet doorgeven?} \\
            }
        }
    }
    & & \\ \\

    \theutterance \stepcounter{utterance}  
    & & & \multicolumn{2}{p{0.3\linewidth}}{
        \cellcolor[rgb]{0.9,0.9,0.9}{
            \makecell[{{p{\linewidth}}}]{
                \texttt{\tiny{[GM$|$GM]}}
                \texttt{ZEG: Wacht, ik krijg een straf. Mijn P staat op (1,7) en (4,1) bevat A. Ik moet eerst A verplaatsen. Laten we stap voor stap werken. Kun jij jouw huidige posities na je zet doorgeven?} \\
            }
        }
    }
    & & \\ \\

    \theutterance \stepcounter{utterance}  
    & & \multicolumn{4}{p{0.6\linewidth}}{
        \cellcolor[rgb]{0.9,0.9,0.9}{
            \makecell[{{p{\linewidth}}}]{
                \texttt{\tiny{[GM$\rangle$P2]}}
                \texttt{Verplaatst 'P' naar (4,1) met succes. Uw bijgewerkte raster ziet er als volgt uit:} \\
\texttt{```} \\
\texttt{ 1234567} \\
\texttt{╔═╤═════╗} \\
\texttt{║◌│SPD◌◌║ 1} \\
\texttt{╟─┤◌◌◌◌◌║ 2} \\
\texttt{║◌└─┐◌A◌║ 3} \\
\texttt{║◌◌M├─┐◌║ 4} \\
\texttt{║◌T◌│◌│O║ 5} \\
\texttt{╟─┬─┴─┼─╢ 6} \\
\texttt{║◌│◌◌◌│◌║ 7} \\
\texttt{╚═╧═══╧═╝} \\
\texttt{```} \\
\texttt{Je speelt momenteel ronde 7 van maximaal 28.} \\
\texttt{Jullie hebben gezamenlijk 8 van 16 straffen verzameld.} \\
\texttt{De andere speler stuurde dit bericht:} \\
\\ 
\\ 
\texttt{"Wacht, ik krijg een straf. Mijn P staat op (1,7) en (4,1) bevat A. Ik moet eerst A verplaatsen. Laten we stap voor stap werken. Kun jij jouw huidige posities na je zet doorgeven?"} \\
\\ 
\\ 
\texttt{Wat is je volgende opdracht?} \\
            }
        }
    }
    & \\ \\

    \theutterance \stepcounter{utterance}  
    & & & \multicolumn{4}{p{0.6\linewidth}}{
        \cellcolor[rgb]{0.9,0.9,0.9}{
            \makecell[{{p{\linewidth}}}]{
                \texttt{\tiny{[GM$\langle$P2]}}
                \texttt{ZEG: Mijn P is verplaatst naar (4,1). Mijn huidige posities: P(4,1), O(7,5), T(2,5), S(2,1), D(4,1) {-} wacht, D staat ook op (4,1)? Laat me controleren. Nee, D staat op (4,1) niet, P is nu op (4,1). Mijn posities: P(4,1), O(7,5), T(2,5), S(2,1), D(3,1) of (4,1)? Ik zie D op het bord. Mijn huidige posities zijn: P(4,1), O(7,5), T(2,5), S(3,1), D(4,1)? Nee, dat kan niet. Mijn raster toont: rij 1: S op (2,1), P op (4,1), D op (3,1)? Sorry, ik raak in de war. Mijn posities: P(4,1), O(7,5), T(2,5), S(2,1), D(3,1), A(6,3), M(3,4).} \\
            }
        }
    }
     \\ \\

    \theutterance \stepcounter{utterance}  
    & & & \multicolumn{2}{p{0.3\linewidth}}{
        \cellcolor[rgb]{0.9,0.9,0.9}{
            \makecell[{{p{\linewidth}}}]{
                \texttt{\tiny{[GM$|$GM]}}
                \texttt{ZEG: Mijn P is verplaatst naar (4,1). Mijn huidige posities: P(4,1), O(7,5), T(2,5), S(2,1), D(4,1) {-} wacht, D staat ook op (4,1)? Laat me controleren. Nee, D staat op (4,1) niet, P is nu op (4,1). Mijn posities: P(4,1), O(7,5), T(2,5), S(2,1), D(3,1) of (4,1)? Ik zie D op het bord. Mijn huidige posities zijn: P(4,1), O(7,5), T(2,5), S(3,1), D(4,1)? Nee, dat kan niet. Mijn raster toont: rij 1: S op (2,1), P op (4,1), D op (3,1)? Sorry, ik raak in de war. Mijn posities: P(4,1), O(7,5), T(2,5), S(2,1), D(3,1), A(6,3), M(3,4).} \\
            }
        }
    }
    & & \\ \\

    \theutterance \stepcounter{utterance}  
    & & \multicolumn{4}{p{0.6\linewidth}}{
        \cellcolor[rgb]{0.9,0.9,0.9}{
            \makecell[{{p{\linewidth}}}]{
                \texttt{\tiny{[P1$\langle$GM]}}
                \texttt{Je bericht is doorgestuurd naar de andere speler.} \\
\\ 
\\ 
\texttt{Je speelt momenteel ronde 7 van maximaal 28.} \\
\texttt{Jullie hebben gezamenlijk 8 van 16 straffen verzameld.} \\
\texttt{De andere speler stuurde dit bericht:} \\
\\ 
\\ 
\texttt{"Mijn P is verplaatst naar (4,1). Mijn huidige posities: P(4,1), O(7,5), T(2,5), S(2,1), D(4,1) {-} wacht, D staat ook op (4,1)? Laat me controleren. Nee, D staat op (4,1) niet, P is nu op (4,1). Mijn posities: P(4,1), O(7,5), T(2,5), S(2,1), D(3,1) of (4,1)? Ik zie D op het bord. Mijn huidige posities zijn: P(4,1), O(7,5), T(2,5), S(3,1), D(4,1)? Nee, dat kan niet. Mijn raster toont: rij 1: S op (2,1), P op (4,1), D op (3,1)? Sorry, ik raak in de war. Mijn posities: P(4,1), O(7,5), T(2,5), S(2,1), D(3,1), A(6,3), M(3,4)."} \\
\\ 
\\ 
\texttt{Wat is je volgende opdracht?} \\
            }
        }
    }
    & \\ \\

    \theutterance \stepcounter{utterance}  
    & \multicolumn{4}{p{0.6\linewidth}}{
        \cellcolor[rgb]{0.9,0.9,0.9}{
            \makecell[{{p{\linewidth}}}]{
                \texttt{\tiny{[P1$\rangle$GM]}}
                \texttt{ZEG: Geen probleem. Mijn huidige posities zijn: P(1,7), O(1,3), T(3,2), S(3,5), D(4,7), A(3,1), M(5,7). Ons doel is: P(4,1), O(6,1), T(1,2), S(5,2), D(4,3), A(6,4), M(5,5). Ik begin met A verplaatsen van (3,1) naar een lege plek, dan kan P naar (4,1). Oké?} \\
            }
        }
    }
    & & \\ \\

    \theutterance \stepcounter{utterance}  
    & & & \multicolumn{2}{p{0.3\linewidth}}{
        \cellcolor[rgb]{0.9,0.9,0.9}{
            \makecell[{{p{\linewidth}}}]{
                \texttt{\tiny{[GM$|$GM]}}
                \texttt{ZEG: Geen probleem. Mijn huidige posities zijn: P(1,7), O(1,3), T(3,2), S(3,5), D(4,7), A(3,1), M(5,7). Ons doel is: P(4,1), O(6,1), T(1,2), S(5,2), D(4,3), A(6,4), M(5,5). Ik begin met A verplaatsen van (3,1) naar een lege plek, dan kan P naar (4,1). Oké?} \\
            }
        }
    }
    & & \\ \\

    \theutterance \stepcounter{utterance}  
    & & \multicolumn{4}{p{0.6\linewidth}}{
        \cellcolor[rgb]{0.9,0.9,0.9}{
            \makecell[{{p{\linewidth}}}]{
                \texttt{\tiny{[GM$\rangle$P2]}}
                \texttt{Je bericht is doorgestuurd naar de andere speler.} \\
\\ 
\\ 
\texttt{Je speelt momenteel ronde 8 van maximaal 28.} \\
\texttt{Jullie hebben gezamenlijk 8 van 16 straffen verzameld.} \\
\texttt{De andere speler stuurde dit bericht:} \\
\\ 
\\ 
\texttt{"Geen probleem. Mijn huidige posities zijn: P(1,7), O(1,3), T(3,2), S(3,5), D(4,7), A(3,1), M(5,7). Ons doel is: P(4,1), O(6,1), T(1,2), S(5,2), D(4,3), A(6,4), M(5,5). Ik begin met A verplaatsen van (3,1) naar een lege plek, dan kan P naar (4,1). Oké?"} \\
\\ 
\\ 
\texttt{Wat is je volgende opdracht?} \\
            }
        }
    }
    & \\ \\

    \theutterance \stepcounter{utterance}  
    & & & \multicolumn{4}{p{0.6\linewidth}}{
        \cellcolor[rgb]{0.9,0.9,0.9}{
            \makecell[{{p{\linewidth}}}]{
                \texttt{\tiny{[GM$\langle$P2]}}
                \texttt{ZEG: Goed plan! Ik ga ondertussen S verplaatsen van (2,1) naar (5,2). Laat me weten wanneer jij klaar bent met jouw zet.} \\
            }
        }
    }
     \\ \\

    \theutterance \stepcounter{utterance}  
    & & & \multicolumn{2}{p{0.3\linewidth}}{
        \cellcolor[rgb]{0.9,0.9,0.9}{
            \makecell[{{p{\linewidth}}}]{
                \texttt{\tiny{[GM$|$GM]}}
                \texttt{ZEG: Goed plan! Ik ga ondertussen S verplaatsen van (2,1) naar (5,2). Laat me weten wanneer jij klaar bent met jouw zet.} \\
            }
        }
    }
    & & \\ \\

    \theutterance \stepcounter{utterance}  
    & & \multicolumn{4}{p{0.6\linewidth}}{
        \cellcolor[rgb]{0.9,0.9,0.9}{
            \makecell[{{p{\linewidth}}}]{
                \texttt{\tiny{[P1$\langle$GM]}}
                \texttt{Je bericht is doorgestuurd naar de andere speler.} \\
\\ 
\\ 
\texttt{Je speelt momenteel ronde 8 van maximaal 28.} \\
\texttt{Jullie hebben gezamenlijk 8 van 16 straffen verzameld.} \\
\texttt{De andere speler stuurde dit bericht:} \\
\\ 
\\ 
\texttt{"Goed plan! Ik ga ondertussen S verplaatsen van (2,1) naar (5,2). Laat me weten wanneer jij klaar bent met jouw zet."} \\
\\ 
\\ 
\texttt{Wat is je volgende opdracht?} \\
            }
        }
    }
    & \\ \\

    \theutterance \stepcounter{utterance}  
    & \multicolumn{4}{p{0.6\linewidth}}{
        \cellcolor[rgb]{0.9,0.9,0.9}{
            \makecell[{{p{\linewidth}}}]{
                \texttt{\tiny{[P1$\rangle$GM]}}
                \texttt{BEWEEG: A, (6,4)} \\
            }
        }
    }
    & & \\ \\

    \theutterance \stepcounter{utterance}  
    & & & \multicolumn{2}{p{0.3\linewidth}}{
        \cellcolor[rgb]{0.9,0.9,0.9}{
            \makecell[{{p{\linewidth}}}]{
                \texttt{\tiny{[GM$|$GM]}}
                \texttt{BEWEEG: A, (6,4)} \\
            }
        }
    }
    & & \\ \\

    \theutterance \stepcounter{utterance}  
    & & & \multicolumn{2}{p{0.3\linewidth}}{
        \cellcolor[rgb]{0.9,0.9,0.9}{
            \makecell[{{p{\linewidth}}}]{
                \texttt{\tiny{[GM$|$GM]}}
                \texttt{Straf: (6,4) is niet leeg, maar bevat '┐'.} \\
\\ 
\texttt{Jullie hebben gezamenlijk 9 van 16 straffen verzameld.} \\
\texttt{Probeer het opnieuw!} \\
            }
        }
    }
    & & \\ \\

    \theutterance \stepcounter{utterance}  
    & & \multicolumn{4}{p{0.6\linewidth}}{
        \cellcolor[rgb]{0.9,0.9,0.9}{
            \makecell[{{p{\linewidth}}}]{
                \texttt{\tiny{[P1$\langle$GM]}}
                \texttt{Straf: (6,4) is niet leeg, maar bevat '┐'.} \\
\\ 
\texttt{Jullie hebben gezamenlijk 9 van 16 straffen verzameld.} \\
\texttt{Probeer het opnieuw!} \\
            }
        }
    }
    & \\ \\

    \theutterance \stepcounter{utterance}  
    & \multicolumn{4}{p{0.6\linewidth}}{
        \cellcolor[rgb]{0.9,0.9,0.9}{
            \makecell[{{p{\linewidth}}}]{
                \texttt{\tiny{[P1$\rangle$GM]}}
                \texttt{ZEG: Ik krijg weer een straf. (6,4) bevat een rasterlijn op mijn bord, geen ◌. Mijn lege cellen in rij 4 zijn: (1,4), (2,4), (4,4), (5,4), (7,4). Kan A naar (1,4) als tussenstap? Dan kan ik later verder schuiven.} \\
            }
        }
    }
    & & \\ \\

    \theutterance \stepcounter{utterance}  
    & & & \multicolumn{2}{p{0.3\linewidth}}{
        \cellcolor[rgb]{0.9,0.9,0.9}{
            \makecell[{{p{\linewidth}}}]{
                \texttt{\tiny{[GM$|$GM]}}
                \texttt{ZEG: Ik krijg weer een straf. (6,4) bevat een rasterlijn op mijn bord, geen ◌. Mijn lege cellen in rij 4 zijn: (1,4), (2,4), (4,4), (5,4), (7,4). Kan A naar (1,4) als tussenstap? Dan kan ik later verder schuiven.} \\
            }
        }
    }
    & & \\ \\

    \theutterance \stepcounter{utterance}  
    & & \multicolumn{4}{p{0.6\linewidth}}{
        \cellcolor[rgb]{0.9,0.9,0.9}{
            \makecell[{{p{\linewidth}}}]{
                \texttt{\tiny{[GM$\rangle$P2]}}
                \texttt{Je bericht is doorgestuurd naar de andere speler.} \\
\\ 
\\ 
\texttt{Je speelt momenteel ronde 9 van maximaal 28.} \\
\texttt{Jullie hebben gezamenlijk 9 van 16 straffen verzameld.} \\
\texttt{De andere speler stuurde dit bericht:} \\
\\ 
\\ 
\texttt{"Ik krijg weer een straf. (6,4) bevat een rasterlijn op mijn bord, geen ◌. Mijn lege cellen in rij 4 zijn: (1,4), (2,4), (4,4), (5,4), (7,4). Kan A naar (1,4) als tussenstap? Dan kan ik later verder schuiven."} \\
\\ 
\\ 
\texttt{Wat is je volgende opdracht?} \\
            }
        }
    }
    & \\ \\

    \theutterance \stepcounter{utterance}  
    & & & \multicolumn{4}{p{0.6\linewidth}}{
        \cellcolor[rgb]{0.9,0.9,0.9}{
            \makecell[{{p{\linewidth}}}]{
                \texttt{\tiny{[GM$\langle$P2]}}
                \texttt{ZEG: (1,4) is hier leeg, dus gebruik dat als tussenstap. Pas wel op dat we uiteindelijk naar de afgesproken doelen gaan. Misschien moeten we onze doelconfiguratie herzien met alleen posities die op beide borden echt leeg zijn. Wil je dat we opnieuw kijken?} \\
            }
        }
    }
     \\ \\

    \theutterance \stepcounter{utterance}  
    & & & \multicolumn{2}{p{0.3\linewidth}}{
        \cellcolor[rgb]{0.9,0.9,0.9}{
            \makecell[{{p{\linewidth}}}]{
                \texttt{\tiny{[GM$|$GM]}}
                \texttt{ZEG: (1,4) is hier leeg, dus gebruik dat als tussenstap. Pas wel op dat we uiteindelijk naar de afgesproken doelen gaan. Misschien moeten we onze doelconfiguratie herzien met alleen posities die op beide borden echt leeg zijn. Wil je dat we opnieuw kijken?} \\
            }
        }
    }
    & & \\ \\

    \theutterance \stepcounter{utterance}  
    & & \multicolumn{4}{p{0.6\linewidth}}{
        \cellcolor[rgb]{0.9,0.9,0.9}{
            \makecell[{{p{\linewidth}}}]{
                \texttt{\tiny{[P1$\langle$GM]}}
                \texttt{Je bericht is doorgestuurd naar de andere speler.} \\
\\ 
\\ 
\texttt{Je speelt momenteel ronde 9 van maximaal 28.} \\
\texttt{Jullie hebben gezamenlijk 9 van 16 straffen verzameld.} \\
\texttt{De andere speler stuurde dit bericht:} \\
\\ 
\\ 
\texttt{"(1,4) is hier leeg, dus gebruik dat als tussenstap. Pas wel op dat we uiteindelijk naar de afgesproken doelen gaan. Misschien moeten we onze doelconfiguratie herzien met alleen posities die op beide borden echt leeg zijn. Wil je dat we opnieuw kijken?"} \\
\\ 
\\ 
\texttt{Wat is je volgende opdracht?} \\
            }
        }
    }
    & \\ \\

    \theutterance \stepcounter{utterance}  
    & \multicolumn{4}{p{0.6\linewidth}}{
        \cellcolor[rgb]{0.9,0.9,0.9}{
            \makecell[{{p{\linewidth}}}]{
                \texttt{\tiny{[P1$\rangle$GM]}}
                \texttt{BEWEEG: A, (1,4)} \\
            }
        }
    }
    & & \\ \\

    \theutterance \stepcounter{utterance}  
    & & & \multicolumn{2}{p{0.3\linewidth}}{
        \cellcolor[rgb]{0.9,0.9,0.9}{
            \makecell[{{p{\linewidth}}}]{
                \texttt{\tiny{[GM$|$GM]}}
                \texttt{BEWEEG: A, (1,4)} \\
            }
        }
    }
    & & \\ \\

    \theutterance \stepcounter{utterance}  
    & & & \multicolumn{2}{p{0.3\linewidth}}{
        \cellcolor[rgb]{0.9,0.9,0.9}{
            \makecell[{{p{\linewidth}}}]{
                \texttt{\tiny{[GM$|$GM]}}
                \texttt{A moved to (1, 4)} \\
            }
        }
    }
    & & \\ \\

    \theutterance \stepcounter{utterance}  
    & & \multicolumn{4}{p{0.6\linewidth}}{
        \cellcolor[rgb]{0.9,0.9,0.9}{
            \makecell[{{p{\linewidth}}}]{
                \texttt{\tiny{[GM$\rangle$P2]}}
                \texttt{Je bericht is doorgestuurd naar de andere speler.} \\
\\ 
\\ 
\texttt{Je speelt momenteel ronde 10 van maximaal 28.} \\
\texttt{Jullie hebben gezamenlijk 9 van 16 straffen verzameld.} \\
\texttt{De andere speler heeft een voorwerp op zijn rooster verplaatst.} \\
\\ 
\\ 
\texttt{Wat is je volgende opdracht?} \\
            }
        }
    }
    & \\ \\

    \theutterance \stepcounter{utterance}  
    & & & \multicolumn{4}{p{0.6\linewidth}}{
        \cellcolor[rgb]{0.9,0.9,0.9}{
            \makecell[{{p{\linewidth}}}]{
                \texttt{\tiny{[GM$\langle$P2]}}
                \texttt{Ik ga S verplaatsen van (2,1) naar (5,2). Ik controleer: (5,2) is leeg op mijn bord.} \\
\\ 
\texttt{BEWEEG: S, (5,2)} \\
            }
        }
    }
     \\ \\

    \theutterance \stepcounter{utterance}  
    & & & \multicolumn{2}{p{0.3\linewidth}}{
        \cellcolor[rgb]{0.9,0.9,0.9}{
            \makecell[{{p{\linewidth}}}]{
                \texttt{\tiny{[GM$|$GM]}}
                \texttt{Ik ga S verplaatsen van (2,1) naar (5,2). Ik controleer: (5,2) is leeg op mijn bord.} \\
\\ 
\texttt{BEWEEG: S, (5,2)} \\
            }
        }
    }
    & & \\ \\

    \theutterance \stepcounter{utterance}  
    & & & \multicolumn{2}{p{0.3\linewidth}}{
        \cellcolor[rgb]{0.9,0.9,0.9}{
            \makecell[{{p{\linewidth}}}]{
                \texttt{\tiny{[GM$|$GM]}}
                \texttt{Invalid format: head is not empty: 'Ik ga S verplaatsen van (2,1) naar (5,2). Ik controleer: (5,2) is leeg op mijn bord.} \\
\\ 
\texttt{'} \\
            }
        }
    }
    & & \\ \\

    \theutterance \stepcounter{utterance}  
    & & \multicolumn{4}{p{0.6\linewidth}}{
        \cellcolor[rgb]{0.9,0.9,0.9}{
            \makecell[{{p{\linewidth}}}]{
                \texttt{\tiny{[GM$\rangle$P2]}}
                \texttt{Straf: Je bericht mag niets vóór het commando bevatten!} \\
\\ 
\texttt{Zorg ervoor dat je antwoord alleen `ZEG: <BERICHT>` of `BEWEEG: <OBJECT>, (<X>, <Y>)` bevat, en niets anders!} \\
\\ 
\texttt{Jullie hebben gezamenlijk 10 van 16 straffen verzameld.} \\
\texttt{Probeer het opnieuw!} \\
            }
        }
    }
    & \\ \\

    \theutterance \stepcounter{utterance}  
    & & & \multicolumn{4}{p{0.6\linewidth}}{
        \cellcolor[rgb]{0.9,0.9,0.9}{
            \makecell[{{p{\linewidth}}}]{
                \texttt{\tiny{[GM$\langle$P2]}}
                \texttt{BEWEEG: S, (5,2)} \\
            }
        }
    }
     \\ \\

    \theutterance \stepcounter{utterance}  
    & & & \multicolumn{2}{p{0.3\linewidth}}{
        \cellcolor[rgb]{0.9,0.9,0.9}{
            \makecell[{{p{\linewidth}}}]{
                \texttt{\tiny{[GM$|$GM]}}
                \texttt{BEWEEG: S, (5,2)} \\
            }
        }
    }
    & & \\ \\

    \theutterance \stepcounter{utterance}  
    & & & \multicolumn{2}{p{0.3\linewidth}}{
        \cellcolor[rgb]{0.9,0.9,0.9}{
            \makecell[{{p{\linewidth}}}]{
                \texttt{\tiny{[GM$|$GM]}}
                \texttt{S moved to (5, 2)} \\
            }
        }
    }
    & & \\ \\

    \theutterance \stepcounter{utterance}  
    & & \multicolumn{4}{p{0.6\linewidth}}{
        \cellcolor[rgb]{0.9,0.9,0.9}{
            \makecell[{{p{\linewidth}}}]{
                \texttt{\tiny{[P1$\langle$GM]}}
                \texttt{Verplaatst 'A' naar (1,4) met succes. Uw bijgewerkte raster ziet er als volgt uit:} \\
\texttt{```} \\
\texttt{ 1234567} \\
\texttt{╔═╤═════╗} \\
\texttt{║◌│◌◌◌◌◌║ 1} \\
\texttt{╟─┤◌T◌◌◌║ 2} \\
\texttt{║O└─┐◌◌◌║ 3} \\
\texttt{║A◌◌├─┐◌║ 4} \\
\texttt{║◌◌S│◌│◌║ 5} \\
\texttt{╟─┬─┴─┼─╢ 6} \\
\texttt{║P│◌DM│◌║ 7} \\
\texttt{╚═╧═══╧═╝} \\
\texttt{```} \\
\texttt{Je speelt momenteel ronde 10 van maximaal 28.} \\
\texttt{Jullie hebben gezamenlijk 10 van 16 straffen verzameld.} \\
\texttt{De andere speler heeft een voorwerp op zijn rooster verplaatst.} \\
\\ 
\\ 
\texttt{Wat is je volgende opdracht?} \\
            }
        }
    }
    & \\ \\

    \theutterance \stepcounter{utterance}  
    & \multicolumn{4}{p{0.6\linewidth}}{
        \cellcolor[rgb]{0.9,0.9,0.9}{
            \makecell[{{p{\linewidth}}}]{
                \texttt{\tiny{[P1$\rangle$GM]}}
                \texttt{BEWEEG: P, (4,1)} \\
            }
        }
    }
    & & \\ \\

    \theutterance \stepcounter{utterance}  
    & & & \multicolumn{2}{p{0.3\linewidth}}{
        \cellcolor[rgb]{0.9,0.9,0.9}{
            \makecell[{{p{\linewidth}}}]{
                \texttt{\tiny{[GM$|$GM]}}
                \texttt{BEWEEG: P, (4,1)} \\
            }
        }
    }
    & & \\ \\

    \theutterance \stepcounter{utterance}  
    & & & \multicolumn{2}{p{0.3\linewidth}}{
        \cellcolor[rgb]{0.9,0.9,0.9}{
            \makecell[{{p{\linewidth}}}]{
                \texttt{\tiny{[GM$|$GM]}}
                \texttt{P moved to (4, 1)} \\
            }
        }
    }
    & & \\ \\

    \theutterance \stepcounter{utterance}  
    & & \multicolumn{4}{p{0.6\linewidth}}{
        \cellcolor[rgb]{0.9,0.9,0.9}{
            \makecell[{{p{\linewidth}}}]{
                \texttt{\tiny{[GM$\rangle$P2]}}
                \texttt{Verplaatst 'S' naar (5,2) met succes. Uw bijgewerkte raster ziet er als volgt uit:} \\
\texttt{```} \\
\texttt{ 1234567} \\
\texttt{╔═╤═════╗} \\
\texttt{║◌│◌PD◌◌║ 1} \\
\texttt{╟─┤◌◌S◌◌║ 2} \\
\texttt{║◌└─┐◌A◌║ 3} \\
\texttt{║◌◌M├─┐◌║ 4} \\
\texttt{║◌T◌│◌│O║ 5} \\
\texttt{╟─┬─┴─┼─╢ 6} \\
\texttt{║◌│◌◌◌│◌║ 7} \\
\texttt{╚═╧═══╧═╝} \\
\texttt{```} \\
\texttt{Je speelt momenteel ronde 11 van maximaal 28.} \\
\texttt{Jullie hebben gezamenlijk 10 van 16 straffen verzameld.} \\
\texttt{De andere speler heeft een voorwerp op zijn rooster verplaatst.} \\
\\ 
\\ 
\texttt{Wat is je volgende opdracht?} \\
            }
        }
    }
    & \\ \\

    \theutterance \stepcounter{utterance}  
    & & & \multicolumn{4}{p{0.6\linewidth}}{
        \cellcolor[rgb]{0.9,0.9,0.9}{
            \makecell[{{p{\linewidth}}}]{
                \texttt{\tiny{[GM$\langle$P2]}}
                \texttt{Mijn volgende zet: T moet naar (1,2). T staat nu op (2,5). (1,2) is leeg op mijn bord.} \\
\\ 
\texttt{BEWEEG: T, (1,2)} \\
            }
        }
    }
     \\ \\

    \theutterance \stepcounter{utterance}  
    & & & \multicolumn{2}{p{0.3\linewidth}}{
        \cellcolor[rgb]{0.9,0.9,0.9}{
            \makecell[{{p{\linewidth}}}]{
                \texttt{\tiny{[GM$|$GM]}}
                \texttt{Mijn volgende zet: T moet naar (1,2). T staat nu op (2,5). (1,2) is leeg op mijn bord.} \\
\\ 
\texttt{BEWEEG: T, (1,2)} \\
            }
        }
    }
    & & \\ \\

    \theutterance \stepcounter{utterance}  
    & & & \multicolumn{2}{p{0.3\linewidth}}{
        \cellcolor[rgb]{0.9,0.9,0.9}{
            \makecell[{{p{\linewidth}}}]{
                \texttt{\tiny{[GM$|$GM]}}
                \texttt{Invalid format: head is not empty: 'Mijn volgende zet: T moet naar (1,2). T staat nu op (2,5). (1,2) is leeg op mijn bord.} \\
\\ 
\texttt{'} \\
            }
        }
    }
    & & \\ \\

    \theutterance \stepcounter{utterance}  
    & & \multicolumn{4}{p{0.6\linewidth}}{
        \cellcolor[rgb]{0.9,0.9,0.9}{
            \makecell[{{p{\linewidth}}}]{
                \texttt{\tiny{[GM$\rangle$P2]}}
                \texttt{Straf: Je bericht mag niets vóór het commando bevatten!} \\
\\ 
\texttt{Zorg ervoor dat je antwoord alleen `ZEG: <BERICHT>` of `BEWEEG: <OBJECT>, (<X>, <Y>)` bevat, en niets anders!} \\
\\ 
\texttt{Jullie hebben gezamenlijk 11 van 16 straffen verzameld.} \\
\texttt{Probeer het opnieuw!} \\
            }
        }
    }
    & \\ \\

    \theutterance \stepcounter{utterance}  
    & & & \multicolumn{4}{p{0.6\linewidth}}{
        \cellcolor[rgb]{0.9,0.9,0.9}{
            \makecell[{{p{\linewidth}}}]{
                \texttt{\tiny{[GM$\langle$P2]}}
                \texttt{BEWEEG: T, (1,2)} \\
            }
        }
    }
     \\ \\

    \theutterance \stepcounter{utterance}  
    & & & \multicolumn{2}{p{0.3\linewidth}}{
        \cellcolor[rgb]{0.9,0.9,0.9}{
            \makecell[{{p{\linewidth}}}]{
                \texttt{\tiny{[GM$|$GM]}}
                \texttt{BEWEEG: T, (1,2)} \\
            }
        }
    }
    & & \\ \\

    \theutterance \stepcounter{utterance}  
    & & & \multicolumn{2}{p{0.3\linewidth}}{
        \cellcolor[rgb]{0.9,0.9,0.9}{
            \makecell[{{p{\linewidth}}}]{
                \texttt{\tiny{[GM$|$GM]}}
                \texttt{Straf: (1,2) is niet leeg, maar bevat '─'.} \\
\\ 
\texttt{Jullie hebben gezamenlijk 12 van 16 straffen verzameld.} \\
\texttt{Probeer het opnieuw!} \\
            }
        }
    }
    & & \\ \\

    \theutterance \stepcounter{utterance}  
    & & \multicolumn{4}{p{0.6\linewidth}}{
        \cellcolor[rgb]{0.9,0.9,0.9}{
            \makecell[{{p{\linewidth}}}]{
                \texttt{\tiny{[GM$\rangle$P2]}}
                \texttt{Straf: (1,2) is niet leeg, maar bevat '─'.} \\
\\ 
\texttt{Jullie hebben gezamenlijk 12 van 16 straffen verzameld.} \\
\texttt{Probeer het opnieuw!} \\
            }
        }
    }
    & \\ \\

    \theutterance \stepcounter{utterance}  
    & & & \multicolumn{4}{p{0.6\linewidth}}{
        \cellcolor[rgb]{0.9,0.9,0.9}{
            \makecell[{{p{\linewidth}}}]{
                \texttt{\tiny{[GM$\langle$P2]}}
                \texttt{BEWEEG: T, (4,4)} \\
            }
        }
    }
     \\ \\

    \theutterance \stepcounter{utterance}  
    & & & \multicolumn{2}{p{0.3\linewidth}}{
        \cellcolor[rgb]{0.9,0.9,0.9}{
            \makecell[{{p{\linewidth}}}]{
                \texttt{\tiny{[GM$|$GM]}}
                \texttt{BEWEEG: T, (4,4)} \\
            }
        }
    }
    & & \\ \\

    \theutterance \stepcounter{utterance}  
    & & & \multicolumn{2}{p{0.3\linewidth}}{
        \cellcolor[rgb]{0.9,0.9,0.9}{
            \makecell[{{p{\linewidth}}}]{
                \texttt{\tiny{[GM$|$GM]}}
                \texttt{Straf: (4,4) is niet leeg, maar bevat '├'.} \\
\\ 
\texttt{Jullie hebben gezamenlijk 13 van 16 straffen verzameld.} \\
\texttt{Probeer het opnieuw!} \\
            }
        }
    }
    & & \\ \\

    \theutterance \stepcounter{utterance}  
    & & \multicolumn{4}{p{0.6\linewidth}}{
        \cellcolor[rgb]{0.9,0.9,0.9}{
            \makecell[{{p{\linewidth}}}]{
                \texttt{\tiny{[GM$\rangle$P2]}}
                \texttt{Straf: (4,4) is niet leeg, maar bevat '├'.} \\
\\ 
\texttt{Jullie hebben gezamenlijk 13 van 16 straffen verzameld.} \\
\texttt{Probeer het opnieuw!} \\
            }
        }
    }
    & \\ \\

    \theutterance \stepcounter{utterance}  
    & & & \multicolumn{4}{p{0.6\linewidth}}{
        \cellcolor[rgb]{0.9,0.9,0.9}{
            \makecell[{{p{\linewidth}}}]{
                \texttt{\tiny{[GM$\langle$P2]}}
                \texttt{BEWEEG: T, (7,2)} \\
            }
        }
    }
     \\ \\

    \theutterance \stepcounter{utterance}  
    & & & \multicolumn{2}{p{0.3\linewidth}}{
        \cellcolor[rgb]{0.9,0.9,0.9}{
            \makecell[{{p{\linewidth}}}]{
                \texttt{\tiny{[GM$|$GM]}}
                \texttt{BEWEEG: T, (7,2)} \\
            }
        }
    }
    & & \\ \\

    \theutterance \stepcounter{utterance}  
    & & & \multicolumn{2}{p{0.3\linewidth}}{
        \cellcolor[rgb]{0.9,0.9,0.9}{
            \makecell[{{p{\linewidth}}}]{
                \texttt{\tiny{[GM$|$GM]}}
                \texttt{T moved to (7, 2)} \\
            }
        }
    }
    & & \\ \\

    \theutterance \stepcounter{utterance}  
    & & \multicolumn{4}{p{0.6\linewidth}}{
        \cellcolor[rgb]{0.9,0.9,0.9}{
            \makecell[{{p{\linewidth}}}]{
                \texttt{\tiny{[P1$\langle$GM]}}
                \texttt{Verplaatst 'P' naar (4,1) met succes. Uw bijgewerkte raster ziet er als volgt uit:} \\
\texttt{```} \\
\texttt{ 1234567} \\
\texttt{╔═╤═════╗} \\
\texttt{║◌│◌P◌◌◌║ 1} \\
\texttt{╟─┤◌T◌◌◌║ 2} \\
\texttt{║O└─┐◌◌◌║ 3} \\
\texttt{║A◌◌├─┐◌║ 4} \\
\texttt{║◌◌S│◌│◌║ 5} \\
\texttt{╟─┬─┴─┼─╢ 6} \\
\texttt{║◌│◌DM│◌║ 7} \\
\texttt{╚═╧═══╧═╝} \\
\texttt{```} \\
\texttt{Je speelt momenteel ronde 11 van maximaal 28.} \\
\texttt{Jullie hebben gezamenlijk 13 van 16 straffen verzameld.} \\
\texttt{De andere speler heeft een voorwerp op zijn rooster verplaatst.} \\
\\ 
\\ 
\texttt{Wat is je volgende opdracht?} \\
            }
        }
    }
    & \\ \\

    \theutterance \stepcounter{utterance}  
    & \multicolumn{4}{p{0.6\linewidth}}{
        \cellcolor[rgb]{0.9,0.9,0.9}{
            \makecell[{{p{\linewidth}}}]{
                \texttt{\tiny{[P1$\rangle$GM]}}
                \texttt{ZEG: P is nu op (4,1). Mijn huidige posities: P(4,1), O(1,3), T(3,2), S(3,5), D(4,7), A(1,4), M(5,7). Nu ga ik O verplaatsen naar (6,1). Wat is jouw status?} \\
            }
        }
    }
    & & \\ \\

    \theutterance \stepcounter{utterance}  
    & & & \multicolumn{2}{p{0.3\linewidth}}{
        \cellcolor[rgb]{0.9,0.9,0.9}{
            \makecell[{{p{\linewidth}}}]{
                \texttt{\tiny{[GM$|$GM]}}
                \texttt{ZEG: P is nu op (4,1). Mijn huidige posities: P(4,1), O(1,3), T(3,2), S(3,5), D(4,7), A(1,4), M(5,7). Nu ga ik O verplaatsen naar (6,1). Wat is jouw status?} \\
            }
        }
    }
    & & \\ \\

    \theutterance \stepcounter{utterance}  
    & & \multicolumn{4}{p{0.6\linewidth}}{
        \cellcolor[rgb]{0.9,0.9,0.9}{
            \makecell[{{p{\linewidth}}}]{
                \texttt{\tiny{[GM$\rangle$P2]}}
                \texttt{Verplaatst 'T' naar (7,2) met succes. Uw bijgewerkte raster ziet er als volgt uit:} \\
\texttt{```} \\
\texttt{ 1234567} \\
\texttt{╔═╤═════╗} \\
\texttt{║◌│◌PD◌◌║ 1} \\
\texttt{╟─┤◌◌S◌T║ 2} \\
\texttt{║◌└─┐◌A◌║ 3} \\
\texttt{║◌◌M├─┐◌║ 4} \\
\texttt{║◌◌◌│◌│O║ 5} \\
\texttt{╟─┬─┴─┼─╢ 6} \\
\texttt{║◌│◌◌◌│◌║ 7} \\
\texttt{╚═╧═══╧═╝} \\
\texttt{```} \\
\texttt{Je speelt momenteel ronde 12 van maximaal 28.} \\
\texttt{Jullie hebben gezamenlijk 13 van 16 straffen verzameld.} \\
\texttt{De andere speler stuurde dit bericht:} \\
\\ 
\\ 
\texttt{"P is nu op (4,1). Mijn huidige posities: P(4,1), O(1,3), T(3,2), S(3,5), D(4,7), A(1,4), M(5,7). Nu ga ik O verplaatsen naar (6,1). Wat is jouw status?"} \\
\\ 
\\ 
\texttt{Wat is je volgende opdracht?} \\
            }
        }
    }
    & \\ \\

    \theutterance \stepcounter{utterance}  
    & & & \multicolumn{4}{p{0.6\linewidth}}{
        \cellcolor[rgb]{0.9,0.9,0.9}{
            \makecell[{{p{\linewidth}}}]{
                \texttt{\tiny{[GM$\langle$P2]}}
                \texttt{ZEG: Mijn huidige posities: P(4,1), O(7,5), T(7,2), S(5,2), D(3,1), A(6,3), M(3,4). Ik moet voorzichtig zijn met lijnen en kaders, heb al meerdere straffen gehad. Kun jij de exacte lege cellen op jouw bord opnieuw bevestigen? Ik wil veilige doelen kiezen.} \\
            }
        }
    }
     \\ \\

    \theutterance \stepcounter{utterance}  
    & & & \multicolumn{2}{p{0.3\linewidth}}{
        \cellcolor[rgb]{0.9,0.9,0.9}{
            \makecell[{{p{\linewidth}}}]{
                \texttt{\tiny{[GM$|$GM]}}
                \texttt{ZEG: Mijn huidige posities: P(4,1), O(7,5), T(7,2), S(5,2), D(3,1), A(6,3), M(3,4). Ik moet voorzichtig zijn met lijnen en kaders, heb al meerdere straffen gehad. Kun jij de exacte lege cellen op jouw bord opnieuw bevestigen? Ik wil veilige doelen kiezen.} \\
            }
        }
    }
    & & \\ \\

    \theutterance \stepcounter{utterance}  
    & & \multicolumn{4}{p{0.6\linewidth}}{
        \cellcolor[rgb]{0.9,0.9,0.9}{
            \makecell[{{p{\linewidth}}}]{
                \texttt{\tiny{[P1$\langle$GM]}}
                \texttt{Je bericht is doorgestuurd naar de andere speler.} \\
\\ 
\\ 
\texttt{Je speelt momenteel ronde 12 van maximaal 28.} \\
\texttt{Jullie hebben gezamenlijk 13 van 16 straffen verzameld.} \\
\texttt{De andere speler stuurde dit bericht:} \\
\\ 
\\ 
\texttt{"Mijn huidige posities: P(4,1), O(7,5), T(7,2), S(5,2), D(3,1), A(6,3), M(3,4). Ik moet voorzichtig zijn met lijnen en kaders, heb al meerdere straffen gehad. Kun jij de exacte lege cellen op jouw bord opnieuw bevestigen? Ik wil veilige doelen kiezen."} \\
\\ 
\\ 
\texttt{Wat is je volgende opdracht?} \\
            }
        }
    }
    & \\ \\

    \theutterance \stepcounter{utterance}  
    & \multicolumn{4}{p{0.6\linewidth}}{
        \cellcolor[rgb]{0.9,0.9,0.9}{
            \makecell[{{p{\linewidth}}}]{
                \texttt{\tiny{[P1$\rangle$GM]}}
                \texttt{ZEG: Mijn exacte lege cellen (◌) zijn: (1,1),(2,1),(5,1),(6,1),(7,1) in rij1; (1,2),(2,2),(5,2),(6,2),(7,2) in rij2; (2,3),(5,3),(6,3),(7,3) in rij3; (2,4),(3,4),(4,4),(5,4),(7,4) in rij4; (1,5),(2,5),(4,5),(5,5),(6,5),(7,5) in rij5; (1,6),(2,6),(3,6),(4,6),(5,6),(6,6) in rij6; (1,7),(2,7),(3,7),(6,7),(7,7) in rij7. Ons huidige doel was P(4,1)✓, O(6,1), T(1,2), S(5,2), D(4,3), A(6,4), M(5,5). Maar (6,4) en (4,3) zijn op mijn bord niet leeg! Laten we doelen kiezen die in beide lijsten leeg zijn. Zie jij overlap met jouw lege cellen?} \\
            }
        }
    }
    & & \\ \\

    \theutterance \stepcounter{utterance}  
    & & & \multicolumn{2}{p{0.3\linewidth}}{
        \cellcolor[rgb]{0.9,0.9,0.9}{
            \makecell[{{p{\linewidth}}}]{
                \texttt{\tiny{[GM$|$GM]}}
                \texttt{ZEG: Mijn exacte lege cellen (◌) zijn: (1,1),(2,1),(5,1),(6,1),(7,1) in rij1; (1,2),(2,2),(5,2),(6,2),(7,2) in rij2; (2,3),(5,3),(6,3),(7,3) in rij3; (2,4),(3,4),(4,4),(5,4),(7,4) in rij4; (1,5),(2,5),(4,5),(5,5),(6,5),(7,5) in rij5; (1,6),(2,6),(3,6),(4,6),(5,6),(6,6) in rij6; (1,7),(2,7),(3,7),(6,7),(7,7) in rij7. Ons huidige doel was P(4,1)✓, O(6,1), T(1,2), S(5,2), D(4,3), A(6,4), M(5,5). Maar (6,4) en (4,3) zijn op mijn bord niet leeg! Laten we doelen kiezen die in beide lijsten leeg zijn. Zie jij overlap met jouw lege cellen?} \\
            }
        }
    }
    & & \\ \\

    \theutterance \stepcounter{utterance}  
    & & \multicolumn{4}{p{0.6\linewidth}}{
        \cellcolor[rgb]{0.9,0.9,0.9}{
            \makecell[{{p{\linewidth}}}]{
                \texttt{\tiny{[GM$\rangle$P2]}}
                \texttt{Je bericht is doorgestuurd naar de andere speler.} \\
\\ 
\\ 
\texttt{Je speelt momenteel ronde 13 van maximaal 28.} \\
\texttt{Jullie hebben gezamenlijk 13 van 16 straffen verzameld.} \\
\texttt{De andere speler stuurde dit bericht:} \\
\\ 
\\ 
\texttt{"Mijn exacte lege cellen (◌) zijn: (1,1),(2,1),(5,1),(6,1),(7,1) in rij1; (1,2),(2,2),(5,2),(6,2),(7,2) in rij2; (2,3),(5,3),(6,3),(7,3) in rij3; (2,4),(3,4),(4,4),(5,4),(7,4) in rij4; (1,5),(2,5),(4,5),(5,5),(6,5),(7,5) in rij5; (1,6),(2,6),(3,6),(4,6),(5,6),(6,6) in rij6; (1,7),(2,7),(3,7),(6,7),(7,7) in rij7. Ons huidige doel was P(4,1)✓, O(6,1), T(1,2), S(5,2), D(4,3), A(6,4), M(5,5). Maar (6,4) en (4,3) zijn op mijn bord niet leeg! Laten we doelen kiezen die in beide lijsten leeg zijn. Zie jij overlap met jouw lege cellen?"} \\
\\ 
\\ 
\texttt{Wat is je volgende opdracht?} \\
            }
        }
    }
    & \\ \\

    \theutterance \stepcounter{utterance}  
    & & & \multicolumn{4}{p{0.6\linewidth}}{
        \cellcolor[rgb]{0.9,0.9,0.9}{
            \makecell[{{p{\linewidth}}}]{
                \texttt{\tiny{[GM$\langle$P2]}}
                \texttt{Ik vergelijk de lege cellen. Mijn lege cellen waren: (3,1),(5,1),(6,1),(7,1) in rij1; (1,2),(2,2),(3,2),(4,2),(6,2) in rij2 {-} wacht, ik moet opnieuw kijken naar mijn huidige bord. Mijn huidige lege cellen zijn: (1,1) was P maar P is verplaatst, dus (1,1) is nu leeg? Nee, P stond op (1,1) en is naar (4,1). Dus (1,1) is nu leeg. Ook (2,1) is leeg want S stond daar en is naar (5,2). Mijn lege cellen nu: (1,1),(2,1),(5,1),(6,1),(7,1) in rij1; (1,2),(2,2),(3,2),(4,2),(6,2) in rij2; (1,3),(2,3),(3,3),(4,3),(5,3),(7,3) in rij3; (1,4),(2,4),(4,4),(5,4),(6,4),(7,4) in rij4; (1,5),(2,5),(3,5),(4,5),(6,5) in rij5; (1,6),(2,6),(3,6),(4,6),(5,6),(6,6),(7,6) in rij6; (1,7),(2,7),(3,7),(4,7),(5,7),(6,7),(7,7) in rij7.} \\
\\ 
\texttt{De overlap is: (1,1),(2,1),(5,1),(6,1),(7,1); (1,2),(2,2),(6,2); (2,3),(5,3),(6,3),(7,3); (2,4),(4,4),(5,4),(7,4); (1,5),(2,5),(4,5),(5,5),(6,5); (1,6),(2,6),(3,6),(4,6),(5,6),(6,6); (1,7),(2,7),(3,7),(6,7),(7,7).} \\
\\ 
\texttt{Voorstel: P(4,1) al gedaan, O(6,1), T(1,2), S(5,2) al gedaan, D(5,5), A(6,3), M(4,4).} \\
\\ 
\texttt{ZEG: Overlap lege cellen: (1,1),(2,1),(5,1),(6,1),(7,1); (1,2),(2,2),(6,2); (2,3),(5,3),(6,3),(7,3); (2,4),(4,4),(5,4),(7,4); (1,5),(2,5),(4,5),(5,5),(6,5); (1,6),(2,6),(3,6),(4,6),(5,6),(6,6); (1,7),(2,7),(3,7),(6,7),(7,7). Nieuw doelvoorstel: P(4,1)✓, O(6,1), T(1,2), S(5,2)✓, D(5,5), A(6,3), M(4,4). Akkoord?} \\
            }
        }
    }
     \\ \\

    \theutterance \stepcounter{utterance}  
    & & & \multicolumn{2}{p{0.3\linewidth}}{
        \cellcolor[rgb]{0.9,0.9,0.9}{
            \makecell[{{p{\linewidth}}}]{
                \texttt{\tiny{[GM$|$GM]}}
                \texttt{Ik vergelijk de lege cellen. Mijn lege cellen waren: (3,1),(5,1),(6,1),(7,1) in rij1; (1,2),(2,2),(3,2),(4,2),(6,2) in rij2 {-} wacht, ik moet opnieuw kijken naar mijn huidige bord. Mijn huidige lege cellen zijn: (1,1) was P maar P is verplaatst, dus (1,1) is nu leeg? Nee, P stond op (1,1) en is naar (4,1). Dus (1,1) is nu leeg. Ook (2,1) is leeg want S stond daar en is naar (5,2). Mijn lege cellen nu: (1,1),(2,1),(5,1),(6,1),(7,1) in rij1; (1,2),(2,2),(3,2),(4,2),(6,2) in rij2; (1,3),(2,3),(3,3),(4,3),(5,3),(7,3) in rij3; (1,4),(2,4),(4,4),(5,4),(6,4),(7,4) in rij4; (1,5),(2,5),(3,5),(4,5),(6,5) in rij5; (1,6),(2,6),(3,6),(4,6),(5,6),(6,6),(7,6) in rij6; (1,7),(2,7),(3,7),(4,7),(5,7),(6,7),(7,7) in rij7.} \\
\\ 
\texttt{De overlap is: (1,1),(2,1),(5,1),(6,1),(7,1); (1,2),(2,2),(6,2); (2,3),(5,3),(6,3),(7,3); (2,4),(4,4),(5,4),(7,4); (1,5),(2,5),(4,5),(5,5),(6,5); (1,6),(2,6),(3,6),(4,6),(5,6),(6,6); (1,7),(2,7),(3,7),(6,7),(7,7).} \\
\\ 
\texttt{Voorstel: P(4,1) al gedaan, O(6,1), T(1,2), S(5,2) al gedaan, D(5,5), A(6,3), M(4,4).} \\
\\ 
\texttt{ZEG: Overlap lege cellen: (1,1),(2,1),(5,1),(6,1),(7,1); (1,2),(2,2),(6,2); (2,3),(5,3),(6,3),(7,3); (2,4),(4,4),(5,4),(7,4); (1,5),(2,5),(4,5),(5,5),(6,5); (1,6),(2,6),(3,6),(4,6),(5,6),(6,6); (1,7),(2,7),(3,7),(6,7),(7,7). Nieuw doelvoorstel: P(4,1)✓, O(6,1), T(1,2), S(5,2)✓, D(5,5), A(6,3), M(4,4). Akkoord?} \\
            }
        }
    }
    & & \\ \\

    \theutterance \stepcounter{utterance}  
    & & & \multicolumn{2}{p{0.3\linewidth}}{
        \cellcolor[rgb]{0.9,0.9,0.9}{
            \makecell[{{p{\linewidth}}}]{
                \texttt{\tiny{[GM$|$GM]}}
                \texttt{Invalid format: head is not empty: 'Ik vergelijk de lege cellen. Mijn lege cellen waren: (3,1),(5,1),(6,1),(7,1) in rij1; (1,2),(2,2),(3,2),(4,2),(6,2) in rij2 {-} wacht, ik moet opnieuw kijken naar mijn huidige bord. Mijn huidige lege cellen zijn: (1,1) was P maar P is verplaatst, dus (1,1) is nu leeg? Nee, P stond op (1,1) en is naar (4,1). Dus (1,1) is nu leeg. Ook (2,1) is leeg want S stond daar en is naar (5,2). Mijn lege cellen nu: (1,1),(2,1),(5,1),(6,1),(7,1) in rij1; (1,2),(2,2),(3,2),(4,2),(6,2) in rij2; (1,3),(2,3),(3,3),(4,3),(5,3),(7,3) in rij3; (1,4),(2,4),(4,4),(5,4),(6,4),(7,4) in rij4; (1,5),(2,5),(3,5),(4,5),(6,5) in rij5; (1,6),(2,6),(3,6),(4,6),(5,6),(6,6),(7,6) in rij6; (1,7),(2,7),(3,7),(4,7),(5,7),(6,7),(7,7) in rij7.} \\
\\ 
\texttt{De overlap is: (1,1),(2,1),(5,1),(6,1),(7,1); (1,2),(2,2),(6,2); (2,3),(5,3),(6,3),(7,3); (2,4),(4,4),(5,4),(7,4); (1,5),(2,5),(4,5),(5,5),(6,5); (1,6),(2,6),(3,6),(4,6),(5,6),(6,6); (1,7),(2,7),(3,7),(6,7),(7,7).} \\
\\ 
\texttt{Voorstel: P(4,1) al gedaan, O(6,1), T(1,2), S(5,2) al gedaan, D(5,5), A(6,3), M(4,4).} \\
\\ 
\texttt{'} \\
            }
        }
    }
    & & \\ \\

    \theutterance \stepcounter{utterance}  
    & & \multicolumn{4}{p{0.6\linewidth}}{
        \cellcolor[rgb]{0.9,0.9,0.9}{
            \makecell[{{p{\linewidth}}}]{
                \texttt{\tiny{[GM$\rangle$P2]}}
                \texttt{Straf: Je bericht mag niets vóór het commando bevatten!} \\
\\ 
\texttt{Zorg ervoor dat je antwoord alleen `ZEG: <BERICHT>` of `BEWEEG: <OBJECT>, (<X>, <Y>)` bevat, en niets anders!} \\
\\ 
\texttt{Jullie hebben gezamenlijk 14 van 16 straffen verzameld.} \\
\texttt{Probeer het opnieuw!} \\
            }
        }
    }
    & \\ \\

    \theutterance \stepcounter{utterance}  
    & & & \multicolumn{4}{p{0.6\linewidth}}{
        \cellcolor[rgb]{0.9,0.9,0.9}{
            \makecell[{{p{\linewidth}}}]{
                \texttt{\tiny{[GM$\langle$P2]}}
                \texttt{ZEG: Overlap lege cellen: (1,1),(2,1),(5,1),(6,1),(7,1); (1,2),(2,2),(6,2); (2,3),(5,3),(6,3),(7,3); (2,4),(4,4),(5,4),(7,4); (1,5),(2,5),(4,5),(5,5),(6,5); (1,6),(2,6),(3,6),(4,6),(5,6),(6,6); (1,7),(2,7),(3,7),(6,7),(7,7). Nieuw doelvoorstel: P(4,1)ok, O(6,1), T(1,2), S(5,2)ok, D(5,5), A(6,3), M(4,4). Akkoord?} \\
            }
        }
    }
     \\ \\

    \theutterance \stepcounter{utterance}  
    & & & \multicolumn{2}{p{0.3\linewidth}}{
        \cellcolor[rgb]{0.9,0.9,0.9}{
            \makecell[{{p{\linewidth}}}]{
                \texttt{\tiny{[GM$|$GM]}}
                \texttt{ZEG: Overlap lege cellen: (1,1),(2,1),(5,1),(6,1),(7,1); (1,2),(2,2),(6,2); (2,3),(5,3),(6,3),(7,3); (2,4),(4,4),(5,4),(7,4); (1,5),(2,5),(4,5),(5,5),(6,5); (1,6),(2,6),(3,6),(4,6),(5,6),(6,6); (1,7),(2,7),(3,7),(6,7),(7,7). Nieuw doelvoorstel: P(4,1)ok, O(6,1), T(1,2), S(5,2)ok, D(5,5), A(6,3), M(4,4). Akkoord?} \\
            }
        }
    }
    & & \\ \\

    \theutterance \stepcounter{utterance}  
    & & \multicolumn{4}{p{0.6\linewidth}}{
        \cellcolor[rgb]{0.9,0.9,0.9}{
            \makecell[{{p{\linewidth}}}]{
                \texttt{\tiny{[P1$\langle$GM]}}
                \texttt{Je bericht is doorgestuurd naar de andere speler.} \\
\\ 
\\ 
\texttt{Je speelt momenteel ronde 13 van maximaal 28.} \\
\texttt{Jullie hebben gezamenlijk 14 van 16 straffen verzameld.} \\
\texttt{De andere speler stuurde dit bericht:} \\
\\ 
\\ 
\texttt{"Overlap lege cellen: (1,1),(2,1),(5,1),(6,1),(7,1); (1,2),(2,2),(6,2); (2,3),(5,3),(6,3),(7,3); (2,4),(4,4),(5,4),(7,4); (1,5),(2,5),(4,5),(5,5),(6,5); (1,6),(2,6),(3,6),(4,6),(5,6),(6,6); (1,7),(2,7),(3,7),(6,7),(7,7). Nieuw doelvoorstel: P(4,1)ok, O(6,1), T(1,2), S(5,2)ok, D(5,5), A(6,3), M(4,4). Akkoord?"} \\
\\ 
\\ 
\texttt{Wat is je volgende opdracht?} \\
            }
        }
    }
    & \\ \\

    \theutterance \stepcounter{utterance}  
    & \multicolumn{4}{p{0.6\linewidth}}{
        \cellcolor[rgb]{0.9,0.9,0.9}{
            \makecell[{{p{\linewidth}}}]{
                \texttt{\tiny{[P1$\rangle$GM]}}
                \texttt{ZEG: (5,2) is bezet door T op mijn bord, dus S kan daar niet naartoe tenzij T eerst verplaatst. Mijn T staat op (3,2). Het doel (1,2) voor T is leeg op mijn bord. (6,3) is leeg voor A. (4,4) is leeg voor M. (5,5) is leeg voor D. (6,1) is leeg voor O. Laten we deze doelen gebruiken: P(4,1), O(6,1), T(1,2), S(5,2), D(5,5), A(6,3), M(4,4). Maar S(5,2) conflicteert nog met T. Zal ik eerst T naar (1,2) verplaatsen, dan is (5,2) vrij voor S?} \\
            }
        }
    }
    & & \\ \\

    \theutterance \stepcounter{utterance}  
    & & & \multicolumn{2}{p{0.3\linewidth}}{
        \cellcolor[rgb]{0.9,0.9,0.9}{
            \makecell[{{p{\linewidth}}}]{
                \texttt{\tiny{[GM$|$GM]}}
                \texttt{ZEG: (5,2) is bezet door T op mijn bord, dus S kan daar niet naartoe tenzij T eerst verplaatst. Mijn T staat op (3,2). Het doel (1,2) voor T is leeg op mijn bord. (6,3) is leeg voor A. (4,4) is leeg voor M. (5,5) is leeg voor D. (6,1) is leeg voor O. Laten we deze doelen gebruiken: P(4,1), O(6,1), T(1,2), S(5,2), D(5,5), A(6,3), M(4,4). Maar S(5,2) conflicteert nog met T. Zal ik eerst T naar (1,2) verplaatsen, dan is (5,2) vrij voor S?} \\
            }
        }
    }
    & & \\ \\

    \theutterance \stepcounter{utterance}  
    & & \multicolumn{4}{p{0.6\linewidth}}{
        \cellcolor[rgb]{0.9,0.9,0.9}{
            \makecell[{{p{\linewidth}}}]{
                \texttt{\tiny{[GM$\rangle$P2]}}
                \texttt{Je bericht is doorgestuurd naar de andere speler.} \\
\\ 
\\ 
\texttt{Je speelt momenteel ronde 14 van maximaal 28.} \\
\texttt{Jullie hebben gezamenlijk 14 van 16 straffen verzameld.} \\
\texttt{De andere speler stuurde dit bericht:} \\
\\ 
\\ 
\texttt{"(5,2) is bezet door T op mijn bord, dus S kan daar niet naartoe tenzij T eerst verplaatst. Mijn T staat op (3,2). Het doel (1,2) voor T is leeg op mijn bord. (6,3) is leeg voor A. (4,4) is leeg voor M. (5,5) is leeg voor D. (6,1) is leeg voor O. Laten we deze doelen gebruiken: P(4,1), O(6,1), T(1,2), S(5,2), D(5,5), A(6,3), M(4,4). Maar S(5,2) conflicteert nog met T. Zal ik eerst T naar (1,2) verplaatsen, dan is (5,2) vrij voor S?"} \\
\\ 
\\ 
\texttt{Wat is je volgende opdracht?} \\
            }
        }
    }
    & \\ \\

    \theutterance \stepcounter{utterance}  
    & & & \multicolumn{4}{p{0.6\linewidth}}{
        \cellcolor[rgb]{0.9,0.9,0.9}{
            \makecell[{{p{\linewidth}}}]{
                \texttt{\tiny{[GM$\langle$P2]}}
                \texttt{ZEG: Ja, verplaats T eerst naar (1,2). Dan wordt (5,2) vrij. Mijn S staat al op (5,2). Ik ga nu D verplaatsen van (3,1) naar (5,5).} \\
            }
        }
    }
     \\ \\

    \theutterance \stepcounter{utterance}  
    & & & \multicolumn{2}{p{0.3\linewidth}}{
        \cellcolor[rgb]{0.9,0.9,0.9}{
            \makecell[{{p{\linewidth}}}]{
                \texttt{\tiny{[GM$|$GM]}}
                \texttt{ZEG: Ja, verplaats T eerst naar (1,2). Dan wordt (5,2) vrij. Mijn S staat al op (5,2). Ik ga nu D verplaatsen van (3,1) naar (5,5).} \\
            }
        }
    }
    & & \\ \\

    \theutterance \stepcounter{utterance}  
    & & \multicolumn{4}{p{0.6\linewidth}}{
        \cellcolor[rgb]{0.9,0.9,0.9}{
            \makecell[{{p{\linewidth}}}]{
                \texttt{\tiny{[P1$\langle$GM]}}
                \texttt{Je bericht is doorgestuurd naar de andere speler.} \\
\\ 
\\ 
\texttt{Je speelt momenteel ronde 14 van maximaal 28.} \\
\texttt{Jullie hebben gezamenlijk 14 van 16 straffen verzameld.} \\
\texttt{De andere speler stuurde dit bericht:} \\
\\ 
\\ 
\texttt{"Ja, verplaats T eerst naar (1,2). Dan wordt (5,2) vrij. Mijn S staat al op (5,2). Ik ga nu D verplaatsen van (3,1) naar (5,5)."} \\
\\ 
\\ 
\texttt{Wat is je volgende opdracht?} \\
            }
        }
    }
    & \\ \\

    \theutterance \stepcounter{utterance}  
    & \multicolumn{4}{p{0.6\linewidth}}{
        \cellcolor[rgb]{0.9,0.9,0.9}{
            \makecell[{{p{\linewidth}}}]{
                \texttt{\tiny{[P1$\rangle$GM]}}
                \texttt{BEWEEG: T, (1,2)} \\
            }
        }
    }
    & & \\ \\

    \theutterance \stepcounter{utterance}  
    & & & \multicolumn{2}{p{0.3\linewidth}}{
        \cellcolor[rgb]{0.9,0.9,0.9}{
            \makecell[{{p{\linewidth}}}]{
                \texttt{\tiny{[GM$|$GM]}}
                \texttt{BEWEEG: T, (1,2)} \\
            }
        }
    }
    & & \\ \\

    \theutterance \stepcounter{utterance}  
    & & & \multicolumn{2}{p{0.3\linewidth}}{
        \cellcolor[rgb]{0.9,0.9,0.9}{
            \makecell[{{p{\linewidth}}}]{
                \texttt{\tiny{[GM$|$GM]}}
                \texttt{Straf: (1,2) is niet leeg, maar bevat '─'.} \\
\\ 
\texttt{Jullie hebben gezamenlijk 15 van 16 straffen verzameld.} \\
\texttt{Probeer het opnieuw!} \\
            }
        }
    }
    & & \\ \\

    \theutterance \stepcounter{utterance}  
    & & \multicolumn{4}{p{0.6\linewidth}}{
        \cellcolor[rgb]{0.9,0.9,0.9}{
            \makecell[{{p{\linewidth}}}]{
                \texttt{\tiny{[P1$\langle$GM]}}
                \texttt{Straf: (1,2) is niet leeg, maar bevat '─'.} \\
\\ 
\texttt{Jullie hebben gezamenlijk 15 van 16 straffen verzameld.} \\
\texttt{Probeer het opnieuw!} \\
            }
        }
    }
    & \\ \\

    \theutterance \stepcounter{utterance}  
    & \multicolumn{4}{p{0.6\linewidth}}{
        \cellcolor[rgb]{0.9,0.9,0.9}{
            \makecell[{{p{\linewidth}}}]{
                \texttt{\tiny{[P1$\rangle$GM]}}
                \texttt{ZEG: Nee! (1,2) bevat een rasterlijn op mijn bord. Mijn lege cellen in rij 2 zijn: (2,2), (5,2), (6,2), (7,2). Kan T naar (2,2) als doel? Of welke andere lege cel in de overlap lijst werkt voor T?} \\
            }
        }
    }
    & & \\ \\

    \theutterance \stepcounter{utterance}  
    & & & \multicolumn{2}{p{0.3\linewidth}}{
        \cellcolor[rgb]{0.9,0.9,0.9}{
            \makecell[{{p{\linewidth}}}]{
                \texttt{\tiny{[GM$|$GM]}}
                \texttt{ZEG: Nee! (1,2) bevat een rasterlijn op mijn bord. Mijn lege cellen in rij 2 zijn: (2,2), (5,2), (6,2), (7,2). Kan T naar (2,2) als doel? Of welke andere lege cel in de overlap lijst werkt voor T?} \\
            }
        }
    }
    & & \\ \\

    \theutterance \stepcounter{utterance}  
    & & \multicolumn{4}{p{0.6\linewidth}}{
        \cellcolor[rgb]{0.9,0.9,0.9}{
            \makecell[{{p{\linewidth}}}]{
                \texttt{\tiny{[GM$\rangle$P2]}}
                \texttt{Je bericht is doorgestuurd naar de andere speler.} \\
\\ 
\\ 
\texttt{Je speelt momenteel ronde 15 van maximaal 28.} \\
\texttt{Jullie hebben gezamenlijk 15 van 16 straffen verzameld.} \\
\texttt{De andere speler stuurde dit bericht:} \\
\\ 
\\ 
\texttt{"Nee! (1,2) bevat een rasterlijn op mijn bord. Mijn lege cellen in rij 2 zijn: (2,2), (5,2), (6,2), (7,2). Kan T naar (2,2) als doel? Of welke andere lege cel in de overlap lijst werkt voor T?"} \\
\\ 
\\ 
\texttt{Wat is je volgende opdracht?} \\
            }
        }
    }
    & \\ \\

    \theutterance \stepcounter{utterance}  
    & & & \multicolumn{4}{p{0.6\linewidth}}{
        \cellcolor[rgb]{0.9,0.9,0.9}{
            \makecell[{{p{\linewidth}}}]{
                \texttt{\tiny{[GM$\langle$P2]}}
                \texttt{ZEG: (2,2) is in de overlap en leeg op mijn bord. Laten we T(2,2) als doel nemen. Dan kan S(5,2) blijven. Dus nieuwe doelen: P(4,1), O(6,1), T(2,2), S(5,2), D(5,5), A(6,3), M(4,4). Akkoord?} \\
            }
        }
    }
     \\ \\

    \theutterance \stepcounter{utterance}  
    & & & \multicolumn{2}{p{0.3\linewidth}}{
        \cellcolor[rgb]{0.9,0.9,0.9}{
            \makecell[{{p{\linewidth}}}]{
                \texttt{\tiny{[GM$|$GM]}}
                \texttt{ZEG: (2,2) is in de overlap en leeg op mijn bord. Laten we T(2,2) als doel nemen. Dan kan S(5,2) blijven. Dus nieuwe doelen: P(4,1), O(6,1), T(2,2), S(5,2), D(5,5), A(6,3), M(4,4). Akkoord?} \\
            }
        }
    }
    & & \\ \\

    \theutterance \stepcounter{utterance}  
    & & \multicolumn{4}{p{0.6\linewidth}}{
        \cellcolor[rgb]{0.9,0.9,0.9}{
            \makecell[{{p{\linewidth}}}]{
                \texttt{\tiny{[P1$\langle$GM]}}
                \texttt{Je bericht is doorgestuurd naar de andere speler.} \\
\\ 
\\ 
\texttt{Je speelt momenteel ronde 15 van maximaal 28.} \\
\texttt{Jullie hebben gezamenlijk 15 van 16 straffen verzameld.} \\
\texttt{De andere speler stuurde dit bericht:} \\
\\ 
\\ 
\texttt{"(2,2) is in de overlap en leeg op mijn bord. Laten we T(2,2) als doel nemen. Dan kan S(5,2) blijven. Dus nieuwe doelen: P(4,1), O(6,1), T(2,2), S(5,2), D(5,5), A(6,3), M(4,4). Akkoord?"} \\
\\ 
\\ 
\texttt{Wat is je volgende opdracht?} \\
            }
        }
    }
    & \\ \\

    \theutterance \stepcounter{utterance}  
    & \multicolumn{4}{p{0.6\linewidth}}{
        \cellcolor[rgb]{0.9,0.9,0.9}{
            \makecell[{{p{\linewidth}}}]{
                \texttt{\tiny{[P1$\rangle$GM]}}
                \texttt{ZEG: Akkoord! Ik verplaats T naar (2,2).} \\
            }
        }
    }
    & & \\ \\

    \theutterance \stepcounter{utterance}  
    & & & \multicolumn{2}{p{0.3\linewidth}}{
        \cellcolor[rgb]{0.9,0.9,0.9}{
            \makecell[{{p{\linewidth}}}]{
                \texttt{\tiny{[GM$|$GM]}}
                \texttt{ZEG: Akkoord! Ik verplaats T naar (2,2).} \\
            }
        }
    }
    & & \\ \\

    \theutterance \stepcounter{utterance}  
    & & \multicolumn{4}{p{0.6\linewidth}}{
        \cellcolor[rgb]{0.9,0.9,0.9}{
            \makecell[{{p{\linewidth}}}]{
                \texttt{\tiny{[GM$\rangle$P2]}}
                \texttt{Je bericht is doorgestuurd naar de andere speler.} \\
\\ 
\\ 
\texttt{Je speelt momenteel ronde 16 van maximaal 28.} \\
\texttt{Jullie hebben gezamenlijk 15 van 16 straffen verzameld.} \\
\texttt{De andere speler stuurde dit bericht:} \\
\\ 
\\ 
\texttt{"Akkoord! Ik verplaats T naar (2,2)."} \\
\\ 
\\ 
\texttt{Wat is je volgende opdracht?} \\
            }
        }
    }
    & \\ \\

    \theutterance \stepcounter{utterance}  
    & & & \multicolumn{4}{p{0.6\linewidth}}{
        \cellcolor[rgb]{0.9,0.9,0.9}{
            \makecell[{{p{\linewidth}}}]{
                \texttt{\tiny{[GM$\langle$P2]}}
                \texttt{BEWEEG: D, (5,5)} \\
            }
        }
    }
     \\ \\

    \theutterance \stepcounter{utterance}  
    & & & \multicolumn{2}{p{0.3\linewidth}}{
        \cellcolor[rgb]{0.9,0.9,0.9}{
            \makecell[{{p{\linewidth}}}]{
                \texttt{\tiny{[GM$|$GM]}}
                \texttt{BEWEEG: D, (5,5)} \\
            }
        }
    }
    & & \\ \\

    \theutterance \stepcounter{utterance}  
    & & & \multicolumn{2}{p{0.3\linewidth}}{
        \cellcolor[rgb]{0.9,0.9,0.9}{
            \makecell[{{p{\linewidth}}}]{
                \texttt{\tiny{[GM$|$GM]}}
                \texttt{D moved to (5, 5)} \\
            }
        }
    }
    & & \\ \\

    \theutterance \stepcounter{utterance}  
    & & \multicolumn{4}{p{0.6\linewidth}}{
        \cellcolor[rgb]{0.9,0.9,0.9}{
            \makecell[{{p{\linewidth}}}]{
                \texttt{\tiny{[P1$\langle$GM]}}
                \texttt{Je bericht is doorgestuurd naar de andere speler.} \\
\\ 
\\ 
\texttt{Je speelt momenteel ronde 16 van maximaal 28.} \\
\texttt{Jullie hebben gezamenlijk 15 van 16 straffen verzameld.} \\
\texttt{De andere speler heeft een voorwerp op zijn rooster verplaatst.} \\
\\ 
\\ 
\texttt{Wat is je volgende opdracht?} \\
            }
        }
    }
    & \\ \\

    \theutterance \stepcounter{utterance}  
    & \multicolumn{4}{p{0.6\linewidth}}{
        \cellcolor[rgb]{0.9,0.9,0.9}{
            \makecell[{{p{\linewidth}}}]{
                \texttt{\tiny{[P1$\rangle$GM]}}
                \texttt{BEWEEG: T, (2,2)} \\
            }
        }
    }
    & & \\ \\

    \theutterance \stepcounter{utterance}  
    & & & \multicolumn{2}{p{0.3\linewidth}}{
        \cellcolor[rgb]{0.9,0.9,0.9}{
            \makecell[{{p{\linewidth}}}]{
                \texttt{\tiny{[GM$|$GM]}}
                \texttt{BEWEEG: T, (2,2)} \\
            }
        }
    }
    & & \\ \\

    \theutterance \stepcounter{utterance}  
    & & & \multicolumn{2}{p{0.3\linewidth}}{
        \cellcolor[rgb]{0.9,0.9,0.9}{
            \makecell[{{p{\linewidth}}}]{
                \texttt{\tiny{[GM$|$GM]}}
                \texttt{Straf: (2,2) is niet leeg, maar bevat '┤'.} \\
\\ 
\texttt{Jullie hebben gezamenlijk 16 van 16 straffen verzameld.} \\
\texttt{Probeer het opnieuw!} \\
            }
        }
    }
    & & \\ \\

    \theutterance \stepcounter{utterance}  
    & & \multicolumn{4}{p{0.6\linewidth}}{
        \cellcolor[rgb]{0.9,0.9,0.9}{
            \makecell[{{p{\linewidth}}}]{
                \texttt{\tiny{[P1$\langle$GM]}}
                \texttt{Straf: (2,2) is niet leeg, maar bevat '┤'.} \\
\\ 
\texttt{Jullie hebben gezamenlijk 16 van 16 straffen verzameld.} \\
\texttt{Probeer het opnieuw!} \\
            }
        }
    }
    & \\ \\

    \theutterance \stepcounter{utterance}  
    & \multicolumn{4}{p{0.6\linewidth}}{
        \cellcolor[rgb]{0.9,0.9,0.9}{
            \makecell[{{p{\linewidth}}}]{
                \texttt{\tiny{[P1$\rangle$GM]}}
                \texttt{ZEG: We hebben 16 straffen. Het spel is waarschijnlijk voorbij. Sorry, ik had beter moeten opletten met de rasterlijnen. Bedankt voor het spelen.} \\
            }
        }
    }
    & & \\ \\

    \theutterance \stepcounter{utterance}  
    & & & \multicolumn{2}{p{0.3\linewidth}}{
        \cellcolor[rgb]{0.9,0.9,0.9}{
            \makecell[{{p{\linewidth}}}]{
                \texttt{\tiny{[GM$|$GM]}}
                \texttt{ZEG: We hebben 16 straffen. Het spel is waarschijnlijk voorbij. Sorry, ik had beter moeten opletten met de rasterlijnen. Bedankt voor het spelen.} \\
            }
        }
    }
    & & \\ \\

    \theutterance \stepcounter{utterance}  
    & & \multicolumn{4}{p{0.6\linewidth}}{
        \cellcolor[rgb]{0.9,0.9,0.9}{
            \makecell[{{p{\linewidth}}}]{
                \texttt{\tiny{[GM$\rangle$P2]}}
                \texttt{Verplaatst 'D' naar (5,5) met succes. Uw bijgewerkte raster ziet er als volgt uit:} \\
\texttt{```} \\
\texttt{ 1234567} \\
\texttt{╔═╤═════╗} \\
\texttt{║◌│◌P◌◌◌║ 1} \\
\texttt{╟─┤◌◌S◌T║ 2} \\
\texttt{║◌└─┐◌A◌║ 3} \\
\texttt{║◌◌M├─┐◌║ 4} \\
\texttt{║◌◌◌│D│O║ 5} \\
\texttt{╟─┬─┴─┼─╢ 6} \\
\texttt{║◌│◌◌◌│◌║ 7} \\
\texttt{╚═╧═══╧═╝} \\
\texttt{```} \\
\texttt{Je speelt momenteel ronde 17 van maximaal 28.} \\
\texttt{Jullie hebben gezamenlijk 16 van 16 straffen verzameld.} \\
\texttt{De andere speler stuurde dit bericht:} \\
\\ 
\\ 
\texttt{"We hebben 16 straffen. Het spel is waarschijnlijk voorbij. Sorry, ik had beter moeten opletten met de rasterlijnen. Bedankt voor het spelen."} \\
\\ 
\\ 
\texttt{Wat is je volgende opdracht?} \\
            }
        }
    }
    & \\ \\

    \theutterance \stepcounter{utterance}  
    & & & \multicolumn{4}{p{0.6\linewidth}}{
        \cellcolor[rgb]{0.9,0.9,0.9}{
            \makecell[{{p{\linewidth}}}]{
                \texttt{\tiny{[GM$\langle$P2]}}
                \texttt{ZEG: Jammer, we hebben het geprobeerd. Bedankt voor de samenwerking! Wil je nog control of we klaar zijn met het spel?} \\
            }
        }
    }
     \\ \\

    \theutterance \stepcounter{utterance}  
    & & & \multicolumn{2}{p{0.3\linewidth}}{
        \cellcolor[rgb]{0.9,0.9,0.9}{
            \makecell[{{p{\linewidth}}}]{
                \texttt{\tiny{[GM$|$GM]}}
                \texttt{ZEG: Jammer, we hebben het geprobeerd. Bedankt voor de samenwerking! Wil je nog control of we klaar zijn met het spel?} \\
            }
        }
    }
    & & \\ \\

    \theutterance \stepcounter{utterance}  
    & & \multicolumn{4}{p{0.6\linewidth}}{
        \cellcolor[rgb]{0.9,0.9,0.9}{
            \makecell[{{p{\linewidth}}}]{
                \texttt{\tiny{[P1$\langle$GM]}}
                \texttt{Je bericht is doorgestuurd naar de andere speler.} \\
\\ 
\\ 
\texttt{Je speelt momenteel ronde 17 van maximaal 28.} \\
\texttt{Jullie hebben gezamenlijk 16 van 16 straffen verzameld.} \\
\texttt{De andere speler stuurde dit bericht:} \\
\\ 
\\ 
\texttt{"Jammer, we hebben het geprobeerd. Bedankt voor de samenwerking! Wil je nog control of we klaar zijn met het spel?"} \\
\\ 
\\ 
\texttt{Wat is je volgende opdracht?} \\
            }
        }
    }
    & \\ \\

    \theutterance \stepcounter{utterance}  
    & \multicolumn{4}{p{0.6\linewidth}}{
        \cellcolor[rgb]{0.9,0.9,0.9}{
            \makecell[{{p{\linewidth}}}]{
                \texttt{\tiny{[P1$\rangle$GM]}}
                \texttt{ZEG: klaar?} \\
            }
        }
    }
    & & \\ \\

    \theutterance \stepcounter{utterance}  
    & & & \multicolumn{2}{p{0.3\linewidth}}{
        \cellcolor[rgb]{0.9,0.9,0.9}{
            \makecell[{{p{\linewidth}}}]{
                \texttt{\tiny{[GM$|$GM]}}
                \texttt{ZEG: klaar?} \\
            }
        }
    }
    & & \\ \\

    \theutterance \stepcounter{utterance}  
    & & \multicolumn{4}{p{0.6\linewidth}}{
        \cellcolor[rgb]{0.9,0.9,0.9}{
            \makecell[{{p{\linewidth}}}]{
                \texttt{\tiny{[GM$\rangle$P2]}}
                \texttt{Je bericht is doorgestuurd naar de andere speler.} \\
\\ 
\\ 
\texttt{Je speelt momenteel ronde 18 van maximaal 28.} \\
\texttt{Jullie hebben gezamenlijk 16 van 16 straffen verzameld.} \\
\texttt{De andere speler stuurde dit bericht:} \\
\\ 
\\ 
\texttt{"klaar?"} \\
\\ 
\\ 
\texttt{Wat is je volgende opdracht?} \\
            }
        }
    }
    & \\ \\

    \theutterance \stepcounter{utterance}  
    & & & \multicolumn{4}{p{0.6\linewidth}}{
        \cellcolor[rgb]{0.9,0.9,0.9}{
            \makecell[{{p{\linewidth}}}]{
                \texttt{\tiny{[GM$\langle$P2]}}
                \texttt{ZEG: klaar!} \\
            }
        }
    }
     \\ \\

    \theutterance \stepcounter{utterance}  
    & & & \multicolumn{2}{p{0.3\linewidth}}{
        \cellcolor[rgb]{0.9,0.9,0.9}{
            \makecell[{{p{\linewidth}}}]{
                \texttt{\tiny{[GM$|$GM]}}
                \texttt{ZEG: klaar!} \\
            }
        }
    }
    & & \\ \\

    \theutterance \stepcounter{utterance}  
    & & & \multicolumn{2}{p{0.3\linewidth}}{
        \cellcolor[rgb]{0.9,0.9,0.9}{
            \makecell[{{p{\linewidth}}}]{
                \texttt{\tiny{[GM$|$GM]}}
                \texttt{true} \\
            }
        }
    }
    & & \\ \\

    \theutterance \stepcounter{utterance}  
    & & & \multicolumn{2}{p{0.3\linewidth}}{
        \cellcolor[rgb]{0.9,0.9,0.9}{
            \makecell[{{p{\linewidth}}}]{
                \texttt{\tiny{[GM$|$GM]}}
                \texttt{Initial states:} \\
\texttt{```} \\
\texttt{Player 1:} \\
\texttt{ 1234567} \\
\texttt{╔═╤═════╗} \\
\texttt{║◌│◌A◌◌◌║ 1} \\
\texttt{╟─┤◌T◌◌◌║ 2} \\
\texttt{║O└─┐◌◌◌║ 3} \\
\texttt{║◌◌◌├─┐◌║ 4} \\
\texttt{║◌◌S│◌│◌║ 5} \\
\texttt{╟─┬─┴─┼─╢ 6} \\
\texttt{║P│◌DM│◌║ 7} \\
\texttt{╚═╧═══╧═╝} \\
\\ 
\\ 
\texttt{Player 2:} \\
\texttt{ 1234567} \\
\texttt{╔═╤═════╗} \\
\texttt{║P│S◌D◌◌║ 1} \\
\texttt{╟─┤◌◌◌◌◌║ 2} \\
\texttt{║◌└─┐◌A◌║ 3} \\
\texttt{║◌◌M├─┐◌║ 4} \\
\texttt{║◌T◌│◌│O║ 5} \\
\texttt{╟─┬─┴─┼─╢ 6} \\
\texttt{║◌│◌◌◌│◌║ 7} \\
\texttt{╚═╧═══╧═╝} \\
\\ 
\texttt{```} \\
            }
        }
    }
    & & \\ \\

    \theutterance \stepcounter{utterance}  
    & & & \multicolumn{2}{p{0.3\linewidth}}{
        \cellcolor[rgb]{0.9,0.9,0.9}{
            \makecell[{{p{\linewidth}}}]{
                \texttt{\tiny{[GM$|$GM]}}
                \texttt{End states:} \\
\texttt{```} \\
\texttt{Player 1:} \\
\texttt{ 1234567} \\
\texttt{╔═╤═════╗} \\
\texttt{║◌│◌P◌◌◌║ 1} \\
\texttt{╟─┤◌T◌◌◌║ 2} \\
\texttt{║O└─┐◌◌◌║ 3} \\
\texttt{║A◌◌├─┐◌║ 4} \\
\texttt{║◌◌S│◌│◌║ 5} \\
\texttt{╟─┬─┴─┼─╢ 6} \\
\texttt{║◌│◌DM│◌║ 7} \\
\texttt{╚═╧═══╧═╝} \\
\\ 
\\ 
\texttt{Player 2:} \\
\texttt{ 1234567} \\
\texttt{╔═╤═════╗} \\
\texttt{║◌│◌P◌◌◌║ 1} \\
\texttt{╟─┤◌◌S◌T║ 2} \\
\texttt{║◌└─┐◌A◌║ 3} \\
\texttt{║◌◌M├─┐◌║ 4} \\
\texttt{║◌◌◌│D│O║ 5} \\
\texttt{╟─┬─┴─┼─╢ 6} \\
\texttt{║◌│◌◌◌│◌║ 7} \\
\texttt{╚═╧═══╧═╝} \\
\\ 
\texttt{```} \\
            }
        }
    }
    & & \\ \\

    \theutterance \stepcounter{utterance}  
    & & & \multicolumn{2}{p{0.3\linewidth}}{
        \cellcolor[rgb]{0.9,0.9,0.9}{
            \makecell[{{p{\linewidth}}}]{
                \texttt{\tiny{[GM$|$GM]}}
                \texttt{* success: True} \\
\texttt{* lose: False} \\
\texttt{* aborted: False} \\
\texttt{{-}{-}{-}{-}{-}{-}{-}} \\
\texttt{* Shifts: 5.00} \\
\texttt{* Max Shifts: 12.00} \\
\texttt{* Min Shifts: 6.00} \\
\texttt{* End Distance Sum: 23.87} \\
\texttt{* Init Distance Sum: 32.45} \\
\texttt{* Expected Distance Sum: 29.33} \\
\texttt{* Penalties: 16.00} \\
\texttt{* Max Penalties: 16.00} \\
\texttt{* Rounds: 18.00} \\
\texttt{* Max Rounds: 28.00} \\
\texttt{* Object Count: 7.00} \\
            }
        }
    }
    & & \\ \\

\end{supertabular}
}

\end{document}
